 %Este está licenciado sob a Licença Creative Commons Atribuição-CompartilhaIgual 3.0 Não Adaptada. Para ver uma cópia desta licença, visite https://creativecommons.org/licenses/by-sa/3.0/ ou envie uma carta para Creative Commons, PO Box 1866, Mountain View, CA 94042, USA.

%\documentclass[main.tex]{subfiles}
%\begin{document}

\chapter{Problemas de valor inicial}\index{problema de valor inicial}
Neste capítulo, estudaremos metodologias numéricas para aproximar a solução de problemas de valor inicial (PVI) para equações diferenciais ordinárias. Inicialmente, consideraremos problemas escalares de primeira ordem. Depois, mostraremos como as mesmas ideias podem ser estendidas para sistemas de equações diferenciais e para problemas de ordem superior.

Considere um problema de valor inicial de primeira ordem dado por:
\begin{subequations}\label{eq:PVI}
\begin{eqnarray}
  u'(t) &=& f(t,u(t)),~~t>t^{(1)},\label{eq:PVI_EDO}\\
  u(t^{(1)}) &=& a ~~ \text{(condição inicial)}.\label{eq:PVI_CI}
\end{eqnarray}
\end{subequations}

A incógnita é uma função $u(t)$ que satisfaz simultaneamente a equação diferencial \eqref{eq:PVI_EDO} e a condição inicial \eqref{eq:PVI_CI}.

Considere os três exemplos a seguir:
\begin{ex} O seguinte problema é linear não homogêneo:
\begin{equation}
\begin{split}
    u'(t) &=t,\\
    u(0) &= 2.
\end{split}
\end{equation}
\end{ex}

\begin{ex} O seguinte problema é linear homogêneo:
\begin{eqnarray}
   u'(t) &=&u(t),\\
   u(0) &=& 1.
\end{eqnarray}
\end{ex}

\begin{ex} O seguinte problema é não linear e não autônomo:
\begin{eqnarray}
   u'(t) &=&\sin(u(t)^2+\sin(t)),\\
   u(0) &=& a.
\end{eqnarray}
\end{ex}

A solução do primeiro exemplo é $u(t)=t^2/2+2$, enquanto a solução do segundo é $u(t)=e^t$. No terceiro exemplo, uma expressão fechada para a solução não é evidente.

Mesmo quando a existência e a unicidade da solução podem ser garantidas teoricamente, em muitos problemas não é possível expressar a solução em termos de funções elementares. Nesses casos, recorremos a métodos numéricos.

Neste capítulo, estudaremos principalmente métodos que aproximam a solução em um conjunto finito de valores de $t$. Esse conjunto será chamado de \emph{malha} e será denotado por
\begin{equation}
 \{t^{(i)}\}_{i=1}^N=\{t^{(1)},t^{(2)},\ldots,t^{(N)}\}.
\end{equation}
Quando o passo é constante,
\begin{equation}
 t^{(i)}=t^{(1)}+(i-1)h,~~~i=1,2,\ldots,N,
\end{equation}
onde $h>0$ é a distância entre dois pontos consecutivos da malha. Denotaremos por $u^{(i)}$ uma aproximação numérica de $u(t^{(i)})$.

%%%%%%%%%%%%%%%%%%%%
% python
%%%%%%%%%%%%%%%%%%%%
\ifispython
Nos códigos em \verb+Python+ apresentados neste capítulo, assumiremos que as seguintes bibliotecas e módulos estão importados:
\begin{verbatim}
import numpy as np
from numpy import linalg
import matplotlib.pyplot as plt
\end{verbatim}
\fi
%%%%%%%%%%%%%%%%%%%%

\section{Rudimentos da teoria de problemas de valor inicial}
Uma questão fundamental no estudo de problemas de valor inicial consiste em determinar se o problema é \emph{bem posto}. Em particular, queremos saber:
\begin{itemize}
 \item se existe solução;
 \item se a solução é única;
 \item se pequenas perturbações nos dados produzem pequenas perturbações na solução, ao menos em intervalos finitos de tempo.
\end{itemize}

Para formular uma condição suficiente de existência e unicidade, introduzimos o conceito de continuidade de Lipschitz.\index{função!Lipschitz}

\begin{defn}
Sejam $I\subset\mathbb{R}$ e $J\subset\mathbb{R}$ intervalos. Dizemos que $f(t,u)$ é Lipschitz contínua em $u$, uniformemente para $t\in I$, se existe uma constante $L>0$ tal que
\begin{equation}
 |f(t,u)-f(t,v)|\leq L|u-v|,~~~\forall t\in I,~~~\forall u,v\in J.
\end{equation}
\end{defn}

O resultado a seguir fornece uma condição clássica de existência e unicidade local.

\begin{teo}[Teorema de Picard-Lindelöf]\index{teorema!de Picard-Lindelöf}
Considere o problema
\begin{equation}\label{eq:pvi_picard}
  \begin{split}
    u'(t)  &= f(t,u(t)), \\
    u(t^{(1)}) &= a.
  \end{split}
\end{equation}
Suponha que $f$ seja contínua em uma vizinhança de $(t^{(1)},a)$ e Lipschitz contínua em $u$ nessa vizinhança. Então existe $t^{(f)}>t^{(1)}$ tal que o problema admite uma única solução em $[t^{(1)},t^{(f)}]$.
\end{teo}

A hipótese de Lipschitz também permite controlar a propagação de perturbações na condição inicial.

\begin{teo}[Dependência contínua na condição inicial]
Sejam $u(t)$ e $v(t)$ soluções de
\begin{eqnarray}
 u'(t)&=&f(t,u(t)),\\
 v'(t)&=&f(t,v(t)),
\end{eqnarray}
em um mesmo intervalo $[t^{(1)},T]$, e suponha que $f$ seja Lipschitz contínua em $u$ com constante $L$. Então
\begin{equation}
 |u(t)-v(t)|\leq e^{L(t-t^{(1)})}|u(t^{(1)})-v(t^{(1)})|,~~~t\in[t^{(1)},T].
\end{equation}
\end{teo}

Essa estimativa mostra que a solução depende continuamente da condição inicial. Ela também será útil mais adiante para compreender por que erros locais introduzidos por um método numérico podem ser controlados em um intervalo finito.

\subsection*{Exercícios resolvidos}

\begin{exeresol}
A função $f(t,u)=\sqrt{u}$, definida para $u\geq 0$, não é Lipschitz contínua em $u$ em qualquer intervalo que contenha $u=0$, pois
\begin{equation}
\lim_{u\to 0+}\frac{|f(t,u)-f(t,0)|}{|u-0|}
=\lim_{u\to 0+}\frac{\sqrt{u}}{u}
=\lim_{u\to 0+}\frac{1}{\sqrt{u}}
=\infty.
\end{equation}
Mostre que o problema de valor inicial
\begin{eqnarray}
\frac{du}{dt} &=&\sqrt{u},~~u\geq 0,\\
u(0) &=& 0
\end{eqnarray}
não admite solução única para $t\geq 0$.
\end{exeresol}
\begin{resol}
A função identicamente nula, $u(t)=0$, satisfaz a equação diferencial e a condição inicial. Por outro lado, a função\footnote{Esta solução pode ser obtida por separação de variáveis.}
\begin{equation}
u(t)=\frac{t^2}{4},~~~t\geq 0,
\end{equation}
também satisfaz o problema, pois
\begin{equation}
u'(t)=\frac{t}{2}=\sqrt{\frac{t^2}{4}},~~~t\geq 0.
\end{equation}

De fato, para qualquer $t_0\geq 0$, a função
\begin{equation}
u(t)=\left\{
\begin{array}{ll}
 0,&0\leq t\leq t_0,\\
 \frac{(t-t_0)^2}{4},&t>t_0
\end{array}
\right.
\end{equation}
é solução. Portanto, o problema possui infinitas soluções.
\end{resol}

\subsection*{Exercícios}

\construirExer

\section{Método de Euler}\index{método!de Euler}\label{sec:euler}
Nesta seção, construiremos o mais simples dos métodos para resolver problemas de valor inicial: o método de Euler com passo constante.\index{Passo} Por passo constante, queremos dizer que os pontos da malha estão todos igualmente espaçados, isto é:
\begin{equation} t^{(i)}=t^{(1)}+(i-1)h,~~~i=1,2,\ldots,N. \end{equation}
onde $h$ é o passo, isto é, a distância entre dois pontos consecutivos da malha.

Considere então o problema de valor inicial dado por:
\begin{equation}\label{eq:EDO1}
  \begin{split}
    u'(t)  &= f(t,u(t)),~~t>t^{(1)} \\
    u(t^{(1)}) &= a.
  \end{split}
\end{equation}
Ao invés de tentar solucionar o problema para qualquer $t>t^{(1)}$, iremos aproximar $u(t)$ em $t=t^{(2)}$.

Integrando \eqref{eq:EDO1} de $t^{(1)}$ até $t^{(2)}$, obtemos:
\begin{eqnarray}
  \int_{t^{(1)}}^{t^{(2)}} u'(t) \;dt &=& \int_{t^{(1)}}^{t^{(2)}} f(t,u(t)) \; dt\\
  u(t^{(2)})-u(t^{(1)})               &=& \int_{t^{(1)}}^{t^{(2)}} f(t,u(t)) \; dt\\
  u(t^{(2)})                      &=& u(t^{(1)}) +  \int _{t^{(1)}}^{t^{(2)}} f(t,u(t)) \; dt
\end{eqnarray}

Seja $u^{(n)}$ a aproximação de $u(t^{(n)})$. Para obter o método numérico mais simples aproximamos $f$ em $[t^{(1)},t^{(2)}]$ pela função constante $f(t,u(t)) \approx  f(t^{(1)},u^{(1)})$,
\begin{eqnarray}
  u^{(2)} &=&  u^{(1)} +   f(t^{(1)},u^{(1)}) \int _{t^{(1)}}^{t^{(2)}}  \; dt \\
  u^{(2)} &=&  u^{(1)} +   f(t^{(1)},u^{(1)}) (t^{(2)}-t^{(1)}) \\
  u^{(2)} &=&  u^{(1)} + h f(t^{(1)},u^{(1)})
\end{eqnarray}

Este procedimento pode ser repetido para  $t^{(3)}$, $t^{(4)}$, $\ldots$, obtendo, assim, o chamado \emph{método de Euler}:
\begin{equation}\label{eq:euler}
  \begin{split}
    u^{(n+1)}&= u^{(n)} + h\;f(t^{(n)},u^{(n)}),\\
    u^{(1)}&=u(t^{(1)})=a~ \hbox{(condição inicial)}.
  \end{split}
\end{equation}


\begin{ex}
Considere o problema de valor inicial
\begin{equation}
 \begin{split}
  u'(t)&=2u(t)\\
  u(0)&=1
 \end{split}
\end{equation}

cuja solução é $u(t)=e^{2t}$. O método de Euler aplicado a este problema produz o  esquema:
\begin{equation}\label{eq:exemplo_y_2y_euler}
  \begin{split}
    u^{(k+1)}&= u^{(k)}+2hu^{(k)}=(1+2h)u^{(k)}\\
    u^{(1)}&= 1,
  \end{split}
\end{equation}
Suponha que queremos calcular o valor aproximado de $u(1)$ com $h=0,2$. Então os pontos $t^{(1)}=0$, $t^{(2)}=0,2$, $t^{(3)}=0,4$, $t^{(4)}=0,6$, $t^{(5)}=0,8$ e $t^{(6)}=1,0$ formam os seis pontos da malha. As aproximações para a solução nos pontos da malha usando o método de Euler são:
\begin{equation}
\begin{split}
  u(0)  &\approx u^{(1)}=1\\
  u(0,2)&\approx u^{(2)}=(1+2h) u^{(1)}=1,4 u^{(1)}=1,4\\
  u(0,4)&\approx u^{(3)}=1,4 u^{(2)}=1,96\\
  u(0,6)&\approx u^{(4)}=1,4 u^{(3)}=2,744\\
  u(0,8)&\approx u^{(5)}=1,4 u^{(4)}=3,8416\\
  u(1,0)&\approx u^{(6)}=1,4 u^{(5)}=5,37824
  \end{split}
\end{equation}
Essa aproximação é bem grosseira quando comparamos com a solução do problema em $t=1$: $u(1)=e^{2}\approx 7,38906$. Não obstante, se tivéssemos escolhido um passo menor, teríamos obtido uma aproximação melhor. Veja tabela abaixo com valores obtidos com diferentes valores de passo $h$.

\begin{tabular}{|l|l|l|l|l|l|l|l|}%\label{pvi:tab_euler}
\hline
   h&$10^{-1}$&$10^{-2}$&$10^{-3}$&$10^{-4}$&$10^{-5}$&$10^{-6}$&$10^{-7}$\\
   \hline
   $u^{(N)}$& 6,1917364 & 7,2446461 & 7,3743124 & 7,3875786 & 7,3889083 & 7,3890413 & 7,3890546\\
   \hline
\end{tabular}

De fato, podemos mostrar que quando $h$ se aproxima de $0$, a solução aproximada via método de Euler converge para a solução exata $e^2$. Para isto, basta observar que a solução da relação de recorrência \eqref{eq:exemplo_y_2y_euler} é dada por
 \begin{equation} u^{(k)}=(1+2h)^{k-1}. \end{equation}
 Como $t^{(k)}=(k-1)h$ e queremos a solução em $t=1$, temos $(k-1)h=1$. Assim, a solução aproximada pelo método de Euler é dada por:
 \begin{equation}
 u^{(k)}= (1+2h)^{k-1}= (1+2h)^{\frac{1}{h}}.
 \end{equation}
Aplicando o limite $h\to 0+$, temos:
  \begin{equation}
  \lim_{h\to 0+} (1+2h)^{\frac{1}{h}}= e^{2}.
  \end{equation}

%%%%%%%%%%%%%%%%%%%%
% python
%%%%%%%%%%%%%%%%%%%%
\ifispython
Em \verb+Python+, podemos computar a solução numérica deste problema de valor inicial via o método de Euler com o seguinte código:
\begin{verbatim}
#define f(t,u)
def f(t,u):
    return 2*u

#tamanho e num. de passos
h = 0.2
N = 6

#cria vetor t e u
t = np.empty(N)
u = np.copy(t)

#C.I.
t[0] = 0
u[0] = 1

#iteracoes
for i in np.arange(N-1):
    t[i+1] = t[i] + h
    u[i+1] = u[i] + h*f(t[i],u[i])

#imprime
for i,tt in enumerate(t):
    print("%1.1f %1.4f" % (t[i],u[i]))
\end{verbatim}
\fi
%%%%%%%%%%%%%%%%%%%%
\end{ex}

Vamos agora, analisar o desempenho do método de Euler usando um exemplo mais complicado, porém ainda simples suficiente para que possamos obter a solução exata:
\begin{ex}\label{ex:ex_euler_1}
Considere o problema de valor inicial relacionado à equação logística\index{equação!logística}:
\begin{equation}\label{eq:logistica}
  \begin{split}
    u'(t)&= u(t)(1-u(t))\\
    u(0)&= 1/2
  \end{split}
\end{equation}
\end{ex}
Podemos obter a solução exata desta equação usando o método de separação de variáveis\index{método!de separação de variáveis} e o método das frações parciais\index{método das frações parciais}. Para tal escrevemos:
\begin{equation}
\frac{du(t)}{u(t)(1-u(t))}=dt
\end{equation}
O termo $\frac{1}{u(t)(1-u(t))}$ pode ser decomposto em frações parciais como $\frac{1}{u}+\frac{1}{1-u}$ e chegamos na seguinte equação diferencial:
\begin{equation}
\left(\frac{1}{u(t)}+\frac{1}{1-u(t)}\right)du=dt.
\end{equation}
Integrando termo-a-termo, temos a seguinte equação algébrica relacionando $u(t)$ e $t$:
\begin{equation}
\ln(u(t))-\ln\left(1-u(t)\right)=t+C
\end{equation}
Onde $C$ é a constante de integração, que é definida pela condição inicial, isto é, $u=1/2$ em $t=0$. Substituindo, temos $C=0$. O que resulta em:
\begin{equation}
\ln\left(\frac{u(t)}{1-u(t)}\right)=t
\end{equation}
Equivalente a
\begin{equation}
\frac{u(t)}{1-u(t)}=e^{t} \Longrightarrow u(t)=(1-u(t))e^{t} \Longrightarrow (1+e^t)u(t)=e^{t}
\end{equation}
E, finalmente, encontramos a solução exata dada por $u(t)=\frac{e^t}{1+e^{t}}$.

Vejamos, agora, o esquema iterativo produzido pelo método de Euler:
\begin{equation}
 \begin{split}
u^{(k+1)}&= u^{(k)}+h u^{(k)}(1-u^{(k)}), \\
u^{(1)}&= 1/2.
 \end{split}
\end{equation}

\ifisscilab
O seguinte código pode ser usado para implementar no \verb+Scilab+ a recursão acima:
\begin{verbatim}
function u=euler(h,Tmax)
  u= .5;
  itmax = Tmax/h;
  for n=1:itmax
    u= u + h*u*(1-u);
   end
endfunction
\end{verbatim}
o qual pode ser invocado da seguinte forma $\left(h=1e-1, t=2\right)$:
\begin{verbatim}
-->euler(1e-1,2)
 ans  =

    0.8854273
\end{verbatim}

\fi

\ifisoctave
O seguinte código pode ser usado para implementar no \verb+Octave+ a recursão acima:
\begin{verbatim}
function u=euler(h,Tmax)
  u= .5;
  itmax = Tmax/h;
  for n=1:itmax
    u= u + h*u*(1-u);
   end
endfunction
\end{verbatim}
o qual pode ser invocado da seguinte forma $\left(h=1e-1, t=2\right)$:
\begin{verbatim}
>> euler(1e-1,2)
ans =  0.88543
\end{verbatim}
\fi


\ifispython
O seguinte código pode ser usado para implementar no \verb+Python+ a recursão acima:

\begin{verbatim}
def euler(h,Tmax):
	u=.5
  	itmax = int(round(Tmax/h))
	for i in range(itmax):
		u = u + h*u*(1-u)
	return u


h=1e-1
for t in [.5, 1, 2, 3]:
	sol_euler=euler(h,t);
	sol_exata=1/(1+np.exp(-t))
	erro_relativo = np.fabs((sol_euler-sol_exata)/sol_exata)
	print("h=%1.0e - u(%1.1f) =~ %1.7f - erro_relativo = %1.1e" % (h, t, sol_euler, erro_relativo) )

h=1e-2
print()
for t in [.5, 1, 2, 3]:
	sol_euler=euler(h,t);
	sol_exata=1/(1+np.exp(-t))
	erro_relativo = np.fabs((sol_euler-sol_exata)/sol_exata)
	print("h=%1.0e - u(%1.1f) =~ %1.7f - erro_relativo = %1.1e" % (h, t, sol_euler, erro_relativo) )

\end{verbatim}
\fi
Para fins de comparação, calculamos a solução exata e aproximada para alguns valores de $t$ e de passo $h$ e resumimos na tabela abaixo:

%\begin{table}[h]
 % \caption{Tabela comparativa entre método de Euler e solução exata para Problema~\ref{ex_euler_1}.}
 % \label{tab:log}
  \begin{tabular}{|c|c|c|c|}\hline
    $t$ & $\text{Exato}$ & $\text{Euler}~~ h=0,1$ & $\text{Euler}~~ h=0,01$\\\hline
    $0$ & $1/2$ & $0,5$ & $0,5$\\\hline
    $1/2$ & $\frac{e^{1/2}}{1+e^{1/2}}\approx 0,6224593$ & $0,6231476$ & $0,6225316$\\\hline
    $1$ & $\frac{e}{1+e}\approx 0,7310586$ & $0,7334030$ & $0,7312946$\\\hline
    $2$ & $\frac{e^2}{1+e^2}\approx  0,8807971$ & $0,8854273$  & $0,8812533$ \\\hline
    $3$ & $\frac{e^3}{1+e^3}\approx   0,9525741$  & $0,9564754$ & $0,9529609$ \\\hline
  \end{tabular}
%\end{table}

\subsection*{Exercícios Resolvidos}

\begin{exeresol}\label{exeresol:exeresol1} Aproxime a solução do problema de valor inicial
\begin{eqnarray}
     u'(t)&=& -0,5u(t)+2+t\\
            u(0) &=&  8
\end{eqnarray}
Usando os seguintes passos: $h=10^{-1}$, $h=10^{-2}$, $h=10^{-3}$, $h=10^{-4}$ e $h=10^{-5}$ e compare a solução aproximada em $t=1$ com a solução exata dada por:
\begin{equation}
     u(t) = 2t+8e^{-t/2} \Longrightarrow u(1)=2+8e^{-1/2} \approx 6,85224527770107
\end{equation}
\end{exeresol}
\begin{resol} Primeiramente identificamos $f(t,u)=-0,5u+2+t$ e construímos o processo iterativo do método de Euler:
\begin{eqnarray}
  u^{(n+1)}&=&u^{(n)} + h( -0,5u^{(n)}+2+t^{(n)}),~~n=1,2,3,\ldots\\
  u^{(1)}&=&8
\end{eqnarray}

\ifisscilab
O seguinte código pode ser usado para implementar no \verb+Scilab+ a recursão acima:
\begin{verbatim}
function u=euler(h,Tmax)
  u= 8;
  t= 0;
  itmax = Tmax/h;
  for n=1:itmax
    t= (n-1)*h;
    u= u + h*(-0.5*u+2+t);
   end
endfunction
\end{verbatim}
o qual pode ser invocado da seguinte forma:
\begin{verbatim}
-->euler(1e-1,1)
 ans  =

    6.7898955
\end{verbatim}
Podemos construir um vetor com as cinco soluções da seguinte forma:
\begin{verbatim}
 -->S=[euler(1e-1,1) euler(1e-2,1) euler(1e-3,1) euler(1e-4,1) euler(1e-5,1)]
 S  =

    6.7898956   6.846163    6.851639    6.852185    6.8522392
\end{verbatim}

\fi

\ifisoctave
O seguinte código pode ser usado para implementar no \verb+Octave+ a recursão acima:
\begin{verbatim}
function u=euler(h,Tmax)
  u= 8;
  t= 0;
  itmax = Tmax/h;
  for n=1:itmax
    t= (n-1)*h;
    u= u + h*(-0.5*u+2+t);
   end
endfunction
\end{verbatim}
o qual pode ser invocado da seguinte forma:
\begin{verbatim}
>> euler(1e-1,1)
ans =  6.7899
\end{verbatim}
Podemos construir um vetor com as cinco soluções da seguinte forma:
\begin{verbatim}
 >> S=[euler(1e-1,1) euler(1e-2,1) euler(1e-3,1) euler(1e-4,1) euler(1e-5,1)]
S =

   6.7899   6.8462   6.8516   6.8522   6.8522
\end{verbatim}
\fi


\ifispython
O seguinte código pode ser usado para implementar no \verb+Python+ a recursão acima:

\begin{verbatim}
def euler(h,Tmax):
	u=8
  	itmax = int(round(Tmax/h))
	for i in range(itmax):
		t=i*h
		k1 = -0.5*u+2+t
		u = u + h*k1
	return u


sol_exata = 2+8*np.exp(-.5)
h=1e-1
for i in np.arange(1,5):
	sol_euler=euler(h,1);
	erro_relativo = np.fabs((sol_euler-sol_exata)/sol_exata)
	print("h=%1.0e - u(1) =~ %1.7f - erro_relativo = %1.1e" % (h, sol_euler, erro_relativo) )
	h=h/10
\end{verbatim}
\fi
A seguinte tabela resume os resultados obtidos:
\begin{center}
 \begin{tabular}{|l|l|l|l|l|l|l|l|}%\label{pvi:tab_euler}
\hline
   h&$10^{-1}$&$10^{-2}$&$10^{-3}$&$10^{-4}$&$10^{-5}$\\
   \hline
  Euler & 6,7898955 &  6,8461635  &  6,8516386  &  6,8521846  &  6,8522392  \\
   \hline
   $\varepsilon_{rel}$ &9,1e-03 &  8,9e-04  & 8,9e-05&   8,9e-06 &  8,9e-07\\
    \hline
  \end{tabular}
\end{center}


% Veja o gráfico da solução para $h=1, 0.5, 0.1, 0.05$:
% \begin{figure}
% \includegraphics[width=\textwidth]{euler.eps}
% \end{figure}

\end{resol}

\subsection*{Exercícios}

\begin{exer} Resolva o problema de valor inicial a seguir envolvendo uma equação não autônoma\index{equação diferencial!não autônoma}, isto é, quando a função $f(t,u)$ depende explicitamente do tempo. Use passo $h=0,1$ e $h=0,01$. Depois compare com a solução exata dada por $u(t)=2e^{-t}+t-1$ nos instantes $t=0$, $t=1$, $t=2$ e $t=3$.
  \begin{equation}
   \begin{split}
    u'(t)&=-u(t)+t,\\
    u(0)&=1.
   \end{split}
  \end{equation}

\end{exer}
\begin{resp}
% O esquema recursivo de Euler fica dado por:
% \begin{eqnarray}
%   u^{(k+1)}&=&u^{(k)}+h(-u^{(k)}+t^{(k)}),\\
%   u^{(1)}&=&1.
% \end{eqnarray}
\begin{center}
  \begin{tabular}{|c|c|c|c|}\hline
    $t$ &  Exato & Euler~~ $h=0,1$ & Euler~~ $h=0,01$\\\hline
    $0$ &  $1$ & $1$ & $1$\\\hline
    $1$ &   $2e^{-1}\approx 0,7357589$ & $0,6973569$   &   $0,7320647$  \\\hline
    $2$ &   $2e^{-2}+1\approx  1,2706706$ & $ 1,2431533 $   &  $ 1,2679593$     \\\hline
    $3$ &   $2e^{-3}+2\approx 2,0995741$  & $ 2,0847823$ & $2,0980818$   \\\hline
  \end{tabular}
\end{center}
\end{resp}

\begin{exer} Resolva o problema de valor inicial envolvendo uma equação não linear\index{Problema de valor inicial!não linear} usando passo $h=0,1$ e $h=0,01$.
 \begin{equation}
  \begin{split}
    u'(t)&=\cos(u(t)),\\
    u(0)&=0.
  \end{split}
 \end{equation}

Depois compare com a solução exata dada por
\begin{equation}u(t)=\tan^{-1} \left( \frac {e^{2t}-1}{{2 e^t}}
 \right).
\end{equation}
nos instantes $t=0$, $t=1$, $t=2$ e $t=3$.
\end{exer}
\begin{resp}
%  \begin{eqnarray}
%   u^{(k+1)}&=&u^{(k)}+h\cos(u^{(k)})\\
%   u^{(1)}&=&0
% \end{eqnarray}
% Comparação:
\begin{center}
  \begin{tabular}{|c|c|c|c|}\hline
    $t$ &  Exato & Euler~~ $h=0,1$ & Euler~~ $h=0,01$\\\hline
    $0$ &  $0$ & $0$ & $0$\\\hline
    $1$ &   $0,8657695 $ & $ 0.8799602$   &   $0.8671764 $  \\\hline
    $2$ &   $1,3017603 $ & $ 1.3196842 $   &  $  1.3035243$     \\\hline
    $3$ &   $1,4713043 $  & $ 1.4827638 $ & $1.4724512 $   \\\hline
  \end{tabular}
\end{center}
\end{resp}

\begin{exer} Resolva a seguinte problema de valor inicial linear com passo $h=10^{-4}$ via método de Euler e compare a solução obtida com o valor exato $y(t)=e^{\sin(t)}$ em $t=2$:
 \begin{equation}
  \begin{split}
    y'(t)&=\cos(t)y(t)\\
    y(0)&=1.
  \end{split}
 \end{equation}
\end{exer}
\begin{resp}
 Aproximação via Euler: $2,4826529 $, exata: $e^{\sin(2)}\approx 2,4825777 $. Erro relativo aproximado: $ 3\times 10^{-5}$.
\end{resp}

\section{Método de Euler melhorado}\index{Método! de Euler melhorado}\label{sec:sec_euler_mod}
O método de Euler estudado na Seção~\ref{sec:euler} é aplicação bastante restrita devido à sua pequena precisão, isto é, normalmente precisamos escolher um passo $h$ muito pequeno para obter soluções de boa qualidade, o que implica um número elevado de passos e, consequentemente, alto custo computacional.

Nesta seção, construiremos o \emph{método de Euler melhorado} ou \emph{método de Euler modificado} ou, ainda, \emph{método de Heun}. Para tal, considere o problema de valor inicial dado por:
\begin{equation}\label{eq:EDO2}
  \begin{split}
    u'(t)  &= f(t,u(t)),~~t>t^{(1)} \\
    u(t^{(1)}) &= a.
  \end{split}
\end{equation}

Assim como fizemos para o método de Euler, integramos \eqref{eq:EDO2} de $t^{(1)}$ até $t^{(2)}$ e obtemos:
\begin{eqnarray}
  \int_{t^{(1)}}^{t^{(2)}} u'(t) \;dt &=& \int_{t^{(1)}}^{t^{(2)}} f(t,u(t)) \; dt\\
  u(t^{(2)})-u(t^{(1)})               &=& \int_{t^{(1)}}^{t^{(2)}} f(t,u(t)) \; dt\\
  u(t^{(2)})                      &=& u(t^{(1)}) +  \int _{t^{(1)}}^{t^{(2)}} f(t,u(t)) \; dt
\end{eqnarray}
Em vez de aproximar $f(t,u(t))$ como uma constante igual ao seu valor em $t=t^{(1)}$, aplicamos a regra do trapézio (ver \ref{sec:trapezio}) à integral envolvida no lado direito da expressão, isto é:
\begin{equation}\label{eq:euler_mel_eq}
\int _{t^{(1)}}^{t^{(2)}} f(t,u(t)) \; dt = \left[\frac{f\left(t^{(1)},u(t^{(1)})\right)+f\left(t^{(2)},u(t^{(2)})\right)}{2}\right]h + O(h^3)
\end{equation}
onde $h=t^{(2)}-t^{(1)}$.
Como o valor de $u(t^{(2)})$ não é conhecido antes de o passo ser realizado, aproximamos seu valor aplicando o método de Euler:
\begin{eqnarray}
\tilde{u}(t^{(2)})= u(t^{(1)})+h f\left(t^{(1)},u(t^{(1)})\right)
\end{eqnarray}
Assim obtemos:
\begin{eqnarray}
  u(t^{(2)})&=& u(t^{(1)}) +  \int _{t^{(1)}}^{t^{(2)}} f(t,u(t)) \; dt\\
  &\approx& u(t^{(1)}) +\left[\frac{f\left(t^{(1)},u(t^{(1)})\right)+f\left(t^{(2)},u(t^{(2)})\right)}{2}\right]h\\
  &\approx& u(t^{(1)}) +\left[\frac{f\left(t^{(1)},u(t^{(1)})\right)+f\left(t^{(2)},\tilde{u}(t^{(2)})\right)}{2}\right]h\\
 % &=& u(t^{(1)}) + \frac{k_1+k_2}{2}h
\end{eqnarray}

Portanto, o método recursivo de Euler melhorado assume a seguinte forma:
\begin{equation}
 \begin{split}
\tilde{u}^{(k+1)}&=u^{(k)}+h f(t^{(k)},u^{(k)}),\\
u^{(k+1)}&=u^{(k)}+\frac{h}{2}\left(f(t^{(k)},u^{(k)})+f(t^{(k+1)},\tilde{u}^{(k+1)})\right),\\
u^{(1)}&=a~~ \hbox{\text{(condição inicial)}}.
 \end{split}
\end{equation}


Que pode ser escrito equivalentemente como:
\begin{equation}
 \begin{split}
k_1&=f(t^{(k)},u^{(k)}),\\
k_2&=f(t^{(k+1)},u^{(k)}+h k_1),\\
u^{(k+1)}&=u^{(k)}+h\dfrac{k_1+k_2}{2},\\
u^{(1)}&=a~~ \hbox{\text{(condição inicial)}}.
 \end{split}
\end{equation}
Aqui $k_1$ e $k_2$ são variáveis auxiliares que representam as inclinações e devem ser calculadas a cada passo. Esta notação é compatível com a notação usada nos métodos de Runge-Kutta, uma família de esquemas iterativos para aproximar problemas de valor inicial, da qual o método de Euler e o método de Euler melhorado são casos particulares. Veremos os métodos de Runge-Kutta na Seção~\ref{sec:sec_RK}.

\subsection*{Exercícios Resolvidos}
\begin{exeresol}\label{exeresol:exeresol1_euler_melhorado} Resolva pelo método de Euler melhorado problema de valor inicial do Exercício Resolvido~\ref{exeresol:exeresol1}:
\begin{eqnarray}
     u'(t)&=& -0,5u(t)+2+t\\
            u(0) &=&  8
\end{eqnarray}

Usando os seguintes passos: $h=10^{-1}$, $h=10^{-2}$, $h=10^{-3}$, $h=10^{-4}$ e $h=10^{-5}$ e compare a solução aproximada em $t=1$ com a solução obtida pelo método de Euler e a solução exata dada por:
\begin{equation}
     u(t) = 2t+8e^{-t/2} \Longrightarrow u(1)=2+8e^{-1/2} \approx 6,85224527770107
\end{equation}
\end{exeresol}
\begin{resol} Primeiramente identificamos $f(t,u)=-0,5u+2+t$ e construímos o processo iterativo do método de Euler melhorado:
\begin{eqnarray}
  k_1&=&  f(t^{(n)},u^{(n)})=-0,5u^{(n)}+2+t^{(n)}\\
  \tilde{u}&=& u^{(n)} + hk_1\\
  k_2&=&  f(t^{(n+1)},\tilde{u})=-0,5\tilde{u}+2+t^{(n+1)}\\
  u^{(n+1)}&=&u^{(n)} + \dfrac{h( k_1+k_2)}{2},~~n=1,2,3,\ldots\\
  u^{(1)}&=&8 ~~ \hbox{\text{(condição inicial)}}.
\end{eqnarray}

\ifisscilab
O seguinte código pode ser usado para implementar no \verb+Scilab+ a recursão acima:
\begin{verbatim}
function u=euler_mod(h,Tmax)
  u= 8;
  itmax = Tmax/h;

  for n=1:itmax
    t=(n-1)*h;
    k1 = (-0.5*u + 2 + t);
    u_til = u + h*k1;
    k2 = (-0.5*u_til + 2 + t + h);
    u= u + h * (k1 + k2)/2;
  end
endfunction
\end{verbatim}
o qual pode ser invocado da seguinte forma:
\begin{verbatim}
-->euler_mod(1e-1,1)
 ans  =

    6.8532941
\end{verbatim}
Podemos construir um vetor com as cinco soluções da seguinte forma:
\begin{verbatim}
 -->S=[euler_mod(1e-1,1) euler_mod(1e-2,1) euler_mod(1e-3,1) euler_mod(1e-4,1) euler_mod(1e-5,1)]
S  =

    6.8532949    6.8522554    6.8522454    6.8522453    6.8522453
 \end{verbatim}
\fi
\ifisoctave
O seguinte código pode ser usado para implementar no \verb+Octave+ a recursão acima:
\begin{verbatim}
function u=euler_mod(h,Tmax)
  u= 8;
  itmax = Tmax/h;

  for n=1:itmax
    t=(n-1)*h;
    k1 = (-0.5*u + 2 + t);
    u_til = u + h*k1;
    k2 = (-0.5*u_til + 2 + t + h);
    u= u + h * (k1 + k2)/2;
  end
endfunction
\end{verbatim}
o qual pode ser invocado da seguinte forma:
\begin{verbatim}
>> euler_mod(1e-1,1)
ans =  6.8533
\end{verbatim}
Podemos construir um vetor com as cinco soluções da seguinte forma:
\begin{verbatim}
 >> S=[euler_mod(1e-1,1) euler_mod(1e-2,1) euler_mod(1e-3,1) euler_mod(1e-4,1) euler_mod(1e-5,1)]
S =

    6.8533   6.8523   6.8522   6.8522   6.8522
\end{verbatim}
\fi
\ifispython
O seguinte código pode ser usado para implementar no \verb+Python+ a recursão acima:

\begin{verbatim}
def euler_mod(h,Tmax):
	u=8
  	itmax = int(round(Tmax/h))
	for i in range(itmax):
		t=i*h
		k1 = -0.5*u+2+t
		u_til = u + h*k1
		k2 = -0.5*u_til+2+(t+h)
		u=u+h*(k1+k2)/2
	return u



sol_exata = 2+8*np.exp(-1/2)
for h in [1e-1, 1e-2, 1e-3, 1e-4, 1e-5]:
	sol_euler=euler_mod(h,1);
	erro_relativo = np.fabs((sol_euler-sol_exata)/sol_exata)
	print("h=%1.0e - u(1) =~ %1.7f - erro_relativo = %1.1e" % (h, sol_euler, erro_relativo) )

\end{verbatim}
\fi
A seguinte tabela resume os resultados obtidos:
\begin{center}
 \begin{tabular}{|l|l|l|l|l|l|l|l|}%\label{pvi:tab_euler}
\hline
   h&$10^{-1}$&$10^{-2}$&$10^{-3}$&$10^{-4}$&$10^{-5}$\\
   \hline
   Euler & 6,7898955 &  6,8461635  &  6,8516386  &  6,8521846  &  6,8522392  \\
   \hline
   $\varepsilon_{rel}$ &9,1e-03 &  8,9e-04  & 8,9e-05&   8,9e-06 &  8,9e-07\\
   \hline
  Euler mod. & 6,8532949 &  6,8522554  &  6,8522454  &  6,8522453  &  6,8522453 \\
   \hline
   $\varepsilon_{rel}$ &1,5e-04 &  1,5e-06  & 1,5e-08&   1,5e-10 &  1,5e-12\\
   \hline

   \end{tabular}
\end{center}
% Veja o gráfico da solução para $h=1, 0.5, 0.1, 0.05$:
% \begin{figure}
% \includegraphics[width=\textwidth]{euler.eps}
% \end{figure}
\end{resol}

\subsection*{Exercícios}

\construirExer

\section{Solução de sistemas de equações diferenciais}\index{Sistemas de equações diferenciais}\label{sec:solsistema}

Nas seções \ref{sec:euler} e \ref{sec:sec_euler_mod}, construímos dois métodos numéricos para resolver problemas de valor inicial. Nestas seções, sempre consideramos problemas envolvendo equações diferenciais ordinárias de primeira ordem, isto é:
\begin{eqnarray}
  u'(t)  &=& f(t,u(t)),~~t>t^{(1)} \\
  u(t^{(1)}) &=& a.
\end{eqnarray}
Estas técnicas podem ser diretamente estendidas para resolver numericamente problemas de valor inicial envolvendo sistemas de equações diferenciais ordinárias de primeira ordem, isto é:
\begin{equation}\label{eq:exemplo_sistema}
  \begin{split}
    u_1'(t)  &= f_1(t,u_1(t), u_2(t), u_3(t),\ldots, u_n(t)),~~t>t^{(1)} ,\\
    u_2'(t)  &= f_2(t,u_1(t), u_2(t), u_3(t),\ldots, u_n(t)),~~t>t^{(1)} ,\\
    u_3'(t)  &= f_3(t,u_1(t), u_2(t), u_3(t),\ldots, u_n(t)),~~t>t^{(1)} ,\\
    &\vdots \\
    u_n'(t)  &= f_n(t,u_1(t), u_2(t), u_3(t),\ldots, u_n(t)),~~t>t^{(1)} ,\\
    u_1(t^{(1)}) &= a_1 ,\\
    u_2(t^{(1)}) &= a_2 ,\\
    u_3(t^{(1)}) &= a_3 ,\\
    &\vdots\\
    u_n(t^{(1)}) &= a_n.
  \end{split}
\end{equation}
O Problema~\eqref{eq:exemplo_sistema} pode ser escrito como um problema de primeira ordem envolvendo uma única incógnita, $u(t)$, dada como um vetor de funções $u_j(t)$, isto é:
\begin{equation}u(t)=\left[
\begin{array}{c}
 u_1(t)\\
 u_2(t)\\
 u_3(t)\\
 \vdots\\
 u_n(t)
\end{array}
\right]\end{equation}
De forma que o Problema~\eqref{eq:exemplo_sistema} assuma a seguinte forma:
\begin{eqnarray}
  u'(t)  &=& f(t,u(t)),~~t>t^{(1)} \\
  u(t^{(1)}) &=& a.
\end{eqnarray}
onde
\begin{equation}f(t,u(t))=\left[
\begin{array}{c}
  f_1(t,u_1(t), u_2(t), u_3(t),\ldots, u_n(t))\\
  f_2(t,u_1(t), u_2(t), u_3(t),\ldots, u_n(t))\\
  f_3(t,u_1(t), u_2(t), u_3(t),\ldots, u_n(t))\\
  \vdots\\
  f_n(t,u_1(t), u_2(t), u_3(t),\ldots, u_n(t))\\
  \end{array}
\right]\end{equation}
e
\begin{equation}a=\left[
\begin{array}{c}
 a_1\\
 a_2\\
 a_3\\
 \vdots\\
 a_n
\end{array}
\right]\end{equation}

Veja o  o Exemplo~\ref{eq:eq_exemplo_sistema}
\begin{ex}\label{ex:exemplo_sistema_PVI}Considere o problema de resolver numericamente pelo método de Euler o seguinte sistema de equações diferenciais ordinárias com valores iniciais:
\begin{subequations}\label{eq:eq_exemplo_sistema}
\begin{eqnarray}
x'(t)&=&-y(t),\\
y'(t)&=&x(t),\\
x(0)&=&1,\\
y(0)&=&0.
\end{eqnarray}
\end{subequations}
%cuja solução exata é $x(t)=\cos(t)$ e $y(t)=\sin(t)$.
Para aplicar o método de Euler a este sistema, devemos encarar as duas incógnitas do sistema como entradas de um vetor, ou seja, escrevemos:
 \begin{equation} u(t)=\left[\begin{array}{c}x(t)\\y(t)\end{array}\right]. \end{equation}
 e, portanto, o sistema pode ser escrito como:
\begin{eqnarray}
\left[\begin{array}{c}x^{(k+1)}\\y^{(k+1)}\end{array}\right]=\left[\begin{array}{c}x^{(k)}\\y^{(k)}\end{array}\right]+h\left[\begin{array}{c}-y^{(k)}\\x^{(k)}\end{array}\right].
\end{eqnarray}
Observe que este processo iterativo é equivalente a discretizar as equações do sistema uma a uma, isto é:
\begin{eqnarray}
x^{(k+1)}&=&x^{(k)}-hy^{(k)},\\
y^{(k+1)}&=&y^{(k)}+hx^{(k)},\\
x^{(1)}&=&1,\\
y^{(1)}&=&0,\\
\end{eqnarray}
\end{ex}

\subsection*{Exercícios Resolvidos}
\begin{exeresol} Resolva pelo método de Euler melhorado o seguinte problema de valor inicial para aproximar o valor de $x$ e $y$ entre $t=0$ e $t=1$:
\begin{eqnarray}
x'(t)&=&x(t)-y(t),\\
y'(t)&=&x(t)-y(t)^3,\\
x(0)&=&1.\\
y(0)&=&0.
\end{eqnarray}
\end{exeresol}
\begin{resol}
 Primeiramente, identificamos $u(t)$ como o vetor incógnita:
 \begin{equation} u(t)=\left[\begin{array}{c}x(t)\\y(t)\end{array}\right]. \end{equation}
Depois aplicamos a recursão do método de Euler melhorado dada por:
 \begin{eqnarray}
\tilde{u}^{(k+1)}&=&u^{(k)}+hf(t^{(k)},u^{(k)}),\\
u^{(k+1)}&=&u^{(k)}+\frac{h}{2}\left(f(t^{(k)},u^{(k)})+f(t^{(k+1)},\tilde{u}^{(k+1)})\right),\\
\end{eqnarray}
isto é:
 \begin{eqnarray}
\tilde{x}^{(k+1)}&=&x^{(k)}+h\left(x^{(k)}-y^{(k)}\right)\\
\tilde{y}^{(k+1)}&=&y^{(k)}+h\left(x^{(k)}-{y^{(k)}}^3\right)\\
{x}^{(k+1)}&=&x^{(k)}+\frac{h}{2}\left[\left({x}^{(k)}-{y}^{(k)}\right)+\left(\tilde{x}^{(k+1)}-\tilde{y}^{(k+1)}\right)\right]\\
{y}^{(k+1)}&=&y^{(k)}+\frac{h}{2}\left[\left({x}^{(k)}-{{y}^{(k)}}^3\right)+ \left(\tilde{x}^{(k+1)}-\left.{\tilde{y}^{(k+1)}}\right.^3\right)\right]\\
 \end{eqnarray}

 A tabela a seguir resume os resultados obtidos:


 \begin{center}
 \begin{tabular}{|l|l|l|l|l|l|l|}%\label{pvi:tab_euler}
\hline
   h&&$t=0,2$&$t=0,4$&$t=0,6$&$t=0,8$&$t=1,0$\\
   \hline
   \multirow{2}{*}{$10^{-2}$} &x  &1.1986240 &1.3890564 &1.5654561 &1.7287187 &1.8874532 \\
			       &y&0.2194288 &0.4692676 &0.7206154 &0.9332802 &1.0850012 \\
   \hline
  \multirow{2}{*}{$10^{-3}$} &x  &1.1986201 &1.3890485 &1.5654455 &1.7287066 &1.8874392\\
			       &y &0.2194293 &0.4692707 &0.7206252 &0.9332999 &1.0850259\\
			       \hline


   \multirow{2}{*}{$10^{-4}$} &x  &1.1986201 &1.3890484 &1.5654454 &1.7287065 &1.8874390\\
			       &y&  0.2194293 &0.4692707 &0.7206253 &0.9333001 &1.0850262 \\
   \hline
   \end{tabular}
\end{center}
\ifispython
A seguinte rotina pode ser usada para implementar a solução do sistema:
\begin{verbatim}
def euler_mod(h,Tmax,u1):
  	itmax = int(round(Tmax/h))
	x=np.empty(itmax+1)
	y=np.empty(itmax+1)
	x[0]=u1[0]
	y[0]=u1[1]

	for i in range(itmax):
		t=i*h
		kx1 = (x[i]-y[i])
		ky1 = (x[i]-y[i]**3)

		x_til = x[i] + h*kx1
		y_til = y[i] + h*ky1

		kx2 = (x_til-y_til)
		ky2 = (x_til-y_til**3)

		x[i+1]=x[i]+h*(kx1+kx2)/2
		y[i+1]=y[i]+h*(ky1+ky2)/2

	return [x,y]


Tmax=1 			#tempo maximo de simulacao
u1=np.asarray([1,0])	#condicoes iniciais na forma vetorial
h=1e-4			#passo
sol_euler=euler_mod(h,Tmax,u1);

itmax = int(round(Tmax/h))for t in [0, .2, .4, .6, .8, 1]:
	k=int(round(t/h))
	print("h=%1.0e - x(%1.1f) =~ %1.6f - y(%1.1f) =~ %1.6f" % (h, t, sol_euler[0][k], t, sol_euler[1][k]) )

\end{verbatim}
\fi

\ifisoctave
A seguinte rotina pode ser usada para implementar a solução do sistema:
\begin{verbatim}
function [x,y]=euler_mod(h,Tmax,u1)
  itmax = Tmax/h;
  x=zeros(itmax+1);
  y=zeros(itmax+1);
  x(1)=u1(1);
  y(1)=u1(2);

  for n = 1:itmax
    t=(n-1)*h;
    kx1 = (x(n)-y(n));
    ky1 = (x(n)-y(n)^3);
    x_til = x(n) + h*kx1;
    y_til = y(n) + h*ky1;

    kx2 = (x_til-y_til);
    ky2 = (x_til-y_til^3);

    x(n+1)=x(n)+h*(kx1+kx2)/2;
    y(n+1)=y(n)+h*(ky1+ky2)/2;
  end
endfunction
\end{verbatim}
Que pode ser invocada como:
\begin{verbatim}
h=1e-2
Tmax=1
itmax=Tmax/h
[x,y]=euler_mod(h,Tmax,[1,0]);
[x(1) x(1+itmax*.2) x(1+itmax*.4) x(1+itmax*.6) x(1+itmax*.8) x(1+itmax)]
[y(1) y(1+itmax*.2) y(1+itmax*.4) y(1+itmax*.6) y(1+itmax*.8) y(1+itmax)]
 \end{verbatim}

\fi


\ifisscilab
A seguinte rotina pode ser usada para implementar a solução do sistema:
\begin{verbatim}
function [x,y]=euler_mod(h,Tmax,u1)
  itmax = Tmax/h;
  x=zeros(itmax+1);
  y=zeros(itmax+1);
  x(1)=u1(1);
  y(1)=u1(2);

  for n = 1:itmax
    t=(n-1)*h;
    kx1 = (x(n)-y(n));
    ky1 = (x(n)-y(n)^3);
    x_til = x(n) + h*kx1;
    y_til = y(n) + h*ky1;

    kx2 = (x_til-y_til);
    ky2 = (x_til-y_til^3);

    x(n+1)=x(n)+h*(kx1+kx2)/2;
    y(n+1)=y(n)+h*(ky1+ky2)/2;
  end
endfunction
\end{verbatim}
Que pode ser invocada como:
\begin{verbatim}
h=1e-2
Tmax=1
itmax=Tmax/h
[x,y]=euler_mod(h,Tmax,[1,0]);
[x(1) x(1+itmax*.2) x(1+itmax*.4) x(1+itmax*.6) x(1+itmax*.8) x(1+itmax)]
[y(1) y(1+itmax*.2) y(1+itmax*.4) y(1+itmax*.6) y(1+itmax*.8) y(1+itmax)]
 \end{verbatim}
\fi
\end{resol}

\subsection*{Exercícios}

\construirExer


\section{Solução de equações e sistemas de ordem superior}

Na Seção~\ref{sec:solsistema}, estendemos os métodos de Euler e Euler melhorado visto nas seções \ref{sec:euler} e \ref{sec:sec_euler_mod} para resolver numericamente problemas de valor inicial envolvendo sistemas de equações diferenciais ordinárias de primeira ordem. Nesta seção, estenderemos estas técnicas para resolver alguns tipos de problemas de ordem superior. Para tal, converteremos a equação diferencial em um sistema, incluindo as derivadas da incógnita como novas incógnitas. Vejamos um exemplo:

\begin{ex} Resolva o problema de valor inicial de segunda ordem dado por
\begin{eqnarray}
y''+y'+y&=&\cos(t),\\
y(0)&=&1,\\
y'(0)&=&0,
\end{eqnarray}
A fim de transformar a equação diferencial dada em um sistema de equações de primeira ordem, introduzimos a substituição $w=y'$, de forma que obteremos o sistema:
\begin{eqnarray}
y'&=&w\\
w'&=&-w-y+\cos(t)\\
y(0)&=&1\\
w(0)&=&0
\end{eqnarray}
Este sistema pode ser resolvido usando as técnicas da Seção~\ref{sec:solsistema}.
%Portanto, o método de Euler produz o seguinte processo iterativo:
%\begin{eqnarray}
%y^{(k+1)}&=&y^{(k)}+hw^{(k)},\\
%w^{(k+1)}&=&w^{(k)}-hw^{(k)}-hy^{(k)}+h\cos(t^{(k)}),\\
%y^{(1)}&=&1,\\
%w^{(1)}&=&0.
%\end{eqnarray}
\end{ex}

\subsection*{Exercícios resolvidos}
\begin{exeresol} Considere o seguinte sistema envolvendo uma equação de segunda ordem e uma de primeira ordem:
\begin{eqnarray}
x''(t)-(1- 0,1 z(t))x'(t)+ x(t)&=&0\\
10 z'(t)+z(t)&=&x(t)^2
\end{eqnarray}
sujeito a condições iniciais dadas por:
\begin{eqnarray}
x(0)&=&3\\
x'(0)&=&0\\
z(0)&=&10\\
\end{eqnarray}
Reescreva este sistema como um sistema de três equações de primeira ordem.
\end{exeresol}
\begin{resol}
 Definimos $y(t)=x'(t)$, pelo que o sistema se torna:
\begin{eqnarray}
x'(t)&=&y(t)\\
y'(t)-(1- 0,1 z(t))y(t)+ x(t)&=&0\\
10 z'(t)+z(t)&=&x(t)^2
\end{eqnarray}
 defina o vetor $u(t)$ como:
\begin{equation}u(t)=\left[
\begin{array}{c}
 x(t)\\y(t)\\z(t)
\end{array}
\right]\end{equation}
De forma que:
\begin{eqnarray}
u'(t)&=&\left[
\begin{array}{c}
 x'(t)\\y'(t)\\z'(t)
\end{array}
\right]=\left[
\begin{array}{c}
 y(t)\\
 \left[1- 0,1 z(t)\right]y(t)- x(t)\\
 \left[x(t)^2-z(t)\right]/10
\end{array}
\right]\\
&\hbox{ou}&\\
u'(t)&=&\left[
\begin{array}{c}
 u_1'(t)\\u_2'(t)\\u_3'(t)
\end{array}
\right]=\left[
\begin{array}{c}
 u_2(t)\\
 \left[1- 0,1 u_3(t)\right]u_2(t)- u_1(t)\\
 \left[u_1(t)^2-u_3(t)\right]/10
\end{array}
\right]
\end{eqnarray}
sujeito às condições iniciais dadas por:
\begin{eqnarray}
u(0)&=&\left[
\begin{array}{c}
 x(0)\\y(0)\\z(0)
\end{array}
\right]=\left[
\begin{array}{c}
3\\0\\10
\end{array}
\right]
\end{eqnarray}

\end{resol}


\subsection*{Exercícios}

\begin{exer}Resolva o problema de valor inicial dado por
\begin{eqnarray}
x'&=& -2x + \sqrt{y}\\
y'&=& x - y\\
x(0)&=&0\\
y(0)&=&2\\
\end{eqnarray}
com passo $h=2\cdot 10^{-1}$, $h=2\cdot 10^{-2}$, $h=2\cdot 10^{-3}$ e $h=2\cdot 10^{-4}$ para obter aproximações para $x(2)$ e $y(2)$.
\end{exer}
\begin{resp}

 \begin{center}
 \begin{tabular}{|l|l|l|l|l|l|}%\label{pvi:tab_euler}
\hline
   h                   & &$2\cdot 10^{-1}$&$2\cdot 10^{-2}$&$2\cdot 10^{-3}$&$2\cdot 10^{-4}$\\
   \hline
 \multirow{2}{*}{Euler}&x&0,4302019&0,4355057&0,4358046&0,4358324\\
                       &y& 0,6172935&0,6457760&0,6486383&0,6489245\\
  \hline

 \multirow{2}{*}{Euler mod,}&x&  0,4343130&0,4358269&0,4358354&0,4358355\\
                       &y& 0,6515479&0,6489764&0,6489566&0,6489564 \\

   \hline


 \end{tabular}
\end{center}
\end{resp}

\begin{exer} Considere o problema de segunda ordem dado por:
\begin{eqnarray}
 x''(t)+x'(t)+\sin(x(t))=1,
\end{eqnarray}
sujeito às condições iniciais dadas por:
\begin{eqnarray}
x(0)&=&2,\\
x'(0)&=&0.
\end{eqnarray}
Resolva numericamente para obter o valor de $x(0,5)$, $x(1)$, $x(1,5)$ e $x(2)$ com passo $h=10^{-2}$ e $h=10^{-3}$ via método de Euler modificado.
\end{exer}

\begin{resp}
Introduzindo $y=x'$, obtemos
\begin{eqnarray}
x'&=&y,\\
y'&=&1-y-\sin(x),
\end{eqnarray}
com $x(0)=2$ e $y(0)=0$. Aplicando o método de Euler modificado, obtemos:
\begin{center}
 \begin{tabular}{|l|l|l|l|l|l|}%\label{pvi:tab_euler}
\hline
   h&&$t=0,5$&$t=1,0$&$t=1,5$&$t=2,0$\\
   \hline
   \multirow{2}{*}{$10^{-2}$} &x&2,0097445&2,0344756&2,0704561&2,1165961\\
                              &y&0,0363118&0,0614276&0,0821603&0,1026166\\
   \hline
   \multirow{2}{*}{$10^{-3}$} &x&2,0097441&2,0344753&2,0704560&2,1165963\\
                              &y&0,0363123&0,0614282&0,0821607&0,1026170\\
   \hline
   \end{tabular}
\end{center}
\end{resp}



\section{Erro de truncamento}\index{Ordem! de precisão}\index{Erro! de truncamento}\label{sec:erro_truncamento}
Nas seções \ref{sec:euler} e \ref{sec:sec_euler_mod}, vimos que a redução do passo $h$ melhora, em geral, a aproximação numérica. No Exercício Resolvido~\ref{exeresol:exeresol1_euler_melhorado}, observa-se também uma diferença importante: o erro do método de Euler decresce aproximadamente de forma proporcional a $h$, enquanto o erro do método de Euler melhorado decresce aproximadamente de forma proporcional a $h^2$. Essa observação motiva o conceito de \emph{ordem de precisão}.

\begin{defn}
Considere um método de passo simples escrito na forma
\begin{equation}
u^{(n+1)}=\Phi_h(t^{(n)},u^{(n)}).
\end{equation}
O \emph{erro de truncamento local} no passo $n$ é o erro cometido em um único passo quando o método parte do valor exato:
\begin{equation}
ETL^{(n+1)}=u(t^{(n+1)})-\Phi_h(t^{(n)},u(t^{(n)})).
\end{equation}
Dizemos que o método tem ordem $p$ se
\begin{equation}
ETL^{(n+1)}=O(h^{p+1})
\end{equation}
uniformemente no intervalo considerado.
\end{defn}

\begin{ex}
O método de Euler tem erro de truncamento local $O(h^2)$ e, portanto, é um método de primeira ordem. De fato, pela expansão de Taylor,
\begin{equation}
u(t+h)=u(t)+hu'(t)+\frac{h^2}{2}u''(t)+O(h^3).
\end{equation}
Escolhendo $t=t^{(n)}$ e usando $u'(t^{(n)})=f(t^{(n)},u(t^{(n)}))$, obtemos
\begin{eqnarray}
u(t^{(n+1)})&=&u(t^{(n)})+hf(t^{(n)},u(t^{(n)}))\nonumber\\
&+&\frac{h^2}{2}u''(t^{(n)})+O(h^3).
\end{eqnarray}
O método de Euler retém os dois primeiros termos do lado direito. Logo,
\begin{equation}
ETL^{(n+1)}=\frac{h^2}{2}u''(t^{(n)})+O(h^3)=O(h^2).
\end{equation}
\end{ex}

O erro de truncamento local mede o defeito de um único passo. Na solução numérica completa, os erros introduzidos em passos anteriores são propagados pelos passos seguintes. Por isso, devemos distinguir o erro local do erro global.

\begin{defn}
O \emph{erro global} no ponto $t^{(n)}$ é definido por
\begin{equation}
e^{(n)}=u(t^{(n)})-u^{(n)}.
\end{equation}
Em um intervalo fixo $[t^{(1)},T]$, dizemos que o método converge com ordem $p$ quando
\begin{equation}
\max_{t^{(n)}\leq T}|e^{(n)}|=O(h^p).
\end{equation}
\end{defn}

Para métodos de passo simples estáveis e sob hipóteses usuais de regularidade, um erro local $O(h^{p+1})$ produz erro global $O(h^p)$ em um intervalo fixo. A perda de uma potência de $h$ pode ser compreendida pelo fato de que o número de passos é da ordem de $1/h$. Esse argumento, porém, só é válido quando a propagação dos erros permanece controlada; retornaremos a esse ponto na seção sobre convergência, consistência e estabilidade.

Para o método de Euler, portanto,
\begin{equation}
ETL_{Euler}=O(h^2)~~~\hbox{e}~~~e^{(n)}=O(h).
\end{equation}

\begin{ex}
Vamos analisar o erro de truncamento local do método de Euler melhorado. Partimos da identidade
\begin{eqnarray}
u(t^{(n+1)})&=&u(t^{(n)})+\int_{t^{(n)}}^{t^{(n+1)}}f(t,u(t))\,dt.
\end{eqnarray}
A regra do trapézio fornece
\begin{eqnarray}
\int_{t^{(n)}}^{t^{(n+1)}}f(t,u(t))\,dt
&=&\frac{h}{2}\left[f(t^{(n)},u(t^{(n)}))+f(t^{(n+1)},u(t^{(n+1)}))\right]+O(h^3).
\end{eqnarray}
No método de Euler melhorado, o valor $u(t^{(n+1)})$ dentro da segunda avaliação de $f$ é substituído pelo preditor de Euler
\begin{equation}
\tilde{u}^{(n+1)}=u(t^{(n)})+hf(t^{(n)},u(t^{(n)})).
\end{equation}
Como $\tilde{u}^{(n+1)}-u(t^{(n+1)})=O(h^2)$ e, supondo $f$ suficientemente regular, a alteração correspondente em $f$ também é $O(h^2)$, sua contribuição após a multiplicação por $h$ é $O(h^3)$. Logo,
\begin{equation}
ETL_{Euler~melhorado}=O(h^3),
\end{equation}
e o método tem ordem $p=2$. Sob as hipóteses de estabilidade usuais, seu erro global é $O(h^2)$.
\end{ex}

\subsection*{Exercícios resolvidos}

\construirExeresol

\subsection*{Exercícios}

\construirExer

\begin{exer}Aplique o método de Euler e o método de Euler melhorado para resolver o problema de valor inicial dado por
\begin{eqnarray}
u'&=& -2u + \sqrt{u}\\
u(0)&=&1
\end{eqnarray}
com passo $h=10^{-1}$, $h=10^{-2}$, $h=10^{-3}$, $h=10^{-4}$ e $h=10^{-5}$ para obter aproximações para $u(1)$. Compare com a solução exata dada do problema dada por $u(t) =  \left({1+2 e^{-t}+e^{-2 t}}\right)/{4}$ através do erro relativo e observe a ordem de precisão do método.
\end{exer}
\begin{resp}

\begin{center}
 \begin{tabular}{|l|l|l|l|l|l|l|l|}%\label{pvi:tab_euler}
\hline
   h&$10^{-1}$&$10^{-2}$&$10^{-3}$&$10^{-4}$&$10^{-5}$\\
   \hline
   Euler & 0,4495791 & 0,4660297 & 0,4675999 & 0,4677562 & 0,4677718\\
   \hline
   $\varepsilon_{rel}$ &9,1e-03 &  8,9e-04  & 8,9e-05&   8,9e-06 &  8,9e-07\\
   \hline
  Euler mod, & 0,4686037 & 0,4677811 & 0,4677736 & 0,4677735 & 0,4677735\\
   \hline
   $\varepsilon_{rel}$ & 1,8e-03 & 1,6e-05 & 1,6e-07 & 1,6e-09 & 1,6e-11\\
   \hline
   \end{tabular}
\end{center}
A solução exata vale $u(1)=\frac{1+2e^{-1}+e^{-2}}{4}= \left(\frac{1+e^{-1}}{2}\right)^2\approx 0.467773541395$.
\end{resp}


\begin{exer}Resolva o problema de valor inicial dado por
\begin{eqnarray}
u'&=& \cos(tu(t))\\
u(0)&=&1\\
\end{eqnarray}
com passo $h=10^{-1}$, $h=10^{-2}$, $h=10^{-3}$, $h=10^{-4}$ e $h=10^{-5}$ para obter aproximações para $u(2)$
\end{exer}
\begin{resp}
\begin{center}
 \begin{tabular}{|l|l|l|l|l|l|l|l|}%\label{pvi:tab_euler}
\hline
   h&$10^{-1}$&$10^{-2}$&$10^{-3}$&$10^{-4}$&$10^{-5}$\\
   \hline
   Euler & 1,1617930 & 1,1395726 & 1,1374484 & 1,1372369 & 1,1372157\\
   \hline
  Euler mod & 1,1365230 & 1,1372075 & 1,1372133 & 1,1372134 & 1,1372134\\
   \hline
   \end{tabular}
\end{center}
\end{resp}




\section{Métodos de Runge-Kutta explícitos}\label{sec:sec_RK}\index{método!de Runge-Kutta explícito}
Nas seções anteriores, exploramos os métodos de Euler e Euler modificado para resolver problemas de valor inicial. Neste momento, deve estar claro ao leitor que o método de Euler melhorado produz soluções de melhor qualidade que o método de Euler para a maior parte dos problemas estudados. Isso se deve, conforme Seção~\ref{sec:erro_truncamento}, à ordem de precisão do método.

Os métodos de Runge-Kutta generalizam o esquema do método de Euler melhorado, inserindo mais estágios de cálculo e buscando ordens de precisão mais altas. Nesta seção, trataremos dos métodos de \emph{Runge-Kutta explícitos}\footnote{Existem também os métodos implícitos que serão abordados na Seção~\ref{sec:sec_IRK}. Ver Observação~\ref{obs:obs_imp}.}

Para tal, considere o  problema de valor inicial:
\begin{eqnarray}
  u'(t) &=&f(t,u(t)), \\
  u(t_0) &=&a.
\end{eqnarray}

Integrando a EDO em $[t^{(n)},t^{(n+1)}]$ obtemos
\begin{eqnarray}
  u^{(n+1)}  &=&u^{(n)}  + \int _{t^{(n)}}^{t^{(n+1)}} f(t,u(t)) \; dt
\end{eqnarray}
O método de Euler aproxima a integral no lado direito da expressão acima utilizando apenas o valor de $f(t,u)$ em $t=t^{(n)}$, o método de Euler melhorado aproxima esta integral utilizando os valores de $f(t,u)$ $t=t^{(n)}$ e $t=t^{(n+1)}$. Os métodos de Runge-Kutta explícitos admitem estágios intermediários, utilizando outros valores de $f(t,u)$ nos pontos $\{\tau_1,\tau_2,\ldots,\tau_\nu\}$ dentro do intervalo $[t_n,t^{(n+1)}]$. Veja esquema abaixo:
\begin{center}
\begin{tabular}{c c c c c}
  $u_n$ &      &       &      & $u_{n+1}$ \\
  $|$   &      &       &      &  $|$  \\ \hline
%  $-+-$ & $--$ &  $--$ & $--$ & $-+-$ \\
  $t^{(n)}$ &      &       &      & $t^{(n+1)}$ \\
  $\tau _1$  & $\tau _2$ & $\cdots $ &      & $\tau_\nu $\\
  $t^{(n)}$~~  &~~ $t^{(n)}+c_2h$~~ & $\cdots $ &      & ~~$t^{(n)}+c_\nu h $~~
  \end{tabular}
\end{center}
Observe que $\tau_j=t^{(n)}+c_jh$ com $0=c_1 \leq c_2 \leq \cdots \leq c_\nu \leq 1$, isto é, o primeira ponto sempre coincide com o extremo esquerdo do intervalo, mas o último ponto não precisa ser o extremo direito. Ademais, um mesmo ponto pode ser usado mais de uma vez com aproximações diferentes para $u(\tau_j)$.

Desta forma, aproximamos a integral por um esquema de quadratura com $\nu$ pontos:
\begin{eqnarray}
   \int_{t^{(n)}}^{t^{(n+1)}} f(t,u(t)) \; dt
           &\approx& h\sum_{j=1}^\nu b_j f\left(\tau_j,u(\tau_j)\right)
\end{eqnarray}
onde $b_j$ são os pesos da quadratura numérica que aproxima a integral. Assim como na construção do método de Euler melhorado, não dispomos dos valores $u(\tau_j)$ antes de calculá-los e, por isso, precisamos estimá-los com base nos estágios anteriores:
\begin{equation}\label{eq:RKa}
  \begin{split}
    \tilde{u}_1 &= u^{(n)} \\
    \tilde{u}_2 &= u^{(n)}  + h a_{21}k_1 \\
    \tilde{u}_3 &= u^{(n)}  + h \left[a_{31}k_1 + a_{32}k_2\right] \\
    \tilde{u}_4 &= u^{(n)}  + h \left[a_{41}k_1 + a_{42}k_2 + a_{43}k_3\right] \\
    &\vdots \\
    \tilde{u}_\nu &= u^{(n)}  + h \left[a_{\nu 1}k_1 + a_{\nu 2}k_2 + \cdots +a_{\nu (\nu-1)}k_{\nu-1}\right] \\
    u^{(n+1)}&= u^{(n)} + h [ b_1k_1+b_2k_2+\cdots + b_\nu k_\nu],
  \end{split}
\end{equation}
onde $k_j=f\left(\tau_j,\tilde{u}_{j}\right)$, $A:=\left(a_{ij}\right)$ é a matriz Runge-Kutta (triangular inferior com diagonal zero), $b_j$ são os pesos Runge-Kutta e $c_j$ são os nós Runge-Kutta. Estes coeficientes podem ser escritos de forma compacta em uma tabela conforme a seguir:
\begin{center}
\begin{tabular}{c|c}
  \begin{tabular}{c}\\$c$\\~\\ \end{tabular} &\begin{tabular}{ccc}& $A$&\end{tabular}     \\  \hline
      & $b^T$
\end{tabular}
~~~$=$~~~
\begin{tabular}{c|ccc}
  $c_1$  & $0$      & $0$      & $0$ \\
  $c_2$  & $a_{21}$ & $0$      & $0$ \\
  $c_3$ & $a_{31}$ & $a_{32}$ & $0$\\  \hline
        & $b_1$     & $b_2$     & $b_3$.
\end{tabular}
~~~$=$~~~
\begin{tabular}{c|ccc}
  $c_1$  &          &         &       \\
  $c_2$  & $a_{21}$ &          &      \\
  $c_3$  & $a_{31}$ & $a_{32}$ &      \\  \hline
         & $b_1$    & $b_2$    & $b_3$.
\end{tabular}

Na tabela mais à direita, omitimos os termos obrigatoriamente nulos.
\end{center}

\begin{ex} O método de Euler modificado pode ser escrito conforme:
\begin{eqnarray}
 \tilde{u}_1 &=& u^{(n)}\\
 \tilde{u}_2 &=& u^{(n)} + h k_1\\
 u^{(n+1)} &=& u^{(n)} + h \left[\frac{1}{2}k_1+\frac{1}{2}k_2\right]\\
  \end{eqnarray}
Identificando os coeficientes, obtemos $\nu=2$, $c_1=0$, $c_2=1$, $a_{21}=1$, $b_1=\frac{1}{2}$ e $b_2=\frac{1}{2}$. Escrevendo na forma tabular, temos:
\begin{center}
\begin{tabular}{c|cc}
  \multirow{2}{*}{$c$}    &  \multicolumn{2}{c}{\multirow{2}{*}{$A$}} 	     \\
                          &  \multicolumn{2}{c}{}                            \\
  \hline
                          & \multicolumn{2}{c}{$~~b~~$}
\end{tabular}
~~~=~~~
\begin{tabular}{c|cc}
  $0$ &   &   \\
  $1$ & $1$ &   \\  \hline
      & $\frac{1}{2}$ &$\frac{1}{2}$
\end{tabular}
\end{center}
\end{ex}

 \begin{obs}\label{obs:obs_imp} Nos métodos chamados explícitos, os elementos da diagonal principal (e acima dela) da matriz $A$ devem ser nulos, pois a aproximação em cada estágio é calculada com base nos valores dos estágios anteriores. Nos métodos implícitos, essa restrição é removida, neste caso, o cálculo da aproximação em um estágio envolve a solução de uma equação algébrica.
  \end{obs}

\begin{obs}
Para que um método de Runge-Kutta seja consistente, e portanto tenha ordem pelo menos $1$, é necessário que
\begin{equation}
\sum_{j=1}^{\nu}b_j=1.
\end{equation}
Além disso, na forma usual dos métodos de Runge-Kutta, os nós e os coeficientes dos estágios satisfazem
\begin{equation}
c_i=\sum_{j=1}^{\nu}a_{ij}.
\end{equation}
No caso explícito, $a_{ij}=0$ para $j\geq i$ e $c_1=0$.
\end{obs}


\subsection{Métodos de Runge-Kutta com dois estágios}\label{eq:RK_2_subsec}
Os métodos de Runge-Kutta com dois estágios $\left(\nu=2\right)$ são da seguinte forma:
\begin{eqnarray}
  \tilde{u}_1 &=&u^{(n)} \\
  \tilde{u}_2 &=&u^{(n)}  + h a_{21}k_1 \\
  u^{(n+1)}&=&u^{(n)}  + h [ b_1k_1+b_2k_2],\label{eq:RK2a}
\end{eqnarray}
onde $k_1=f(t^{(n)},u^{(n)})$ e $k_2=f(t^{(n)}+c_2h,\tilde{u}_2)$

Assumindo suavidade suficiente em $f$, usamos o polinômio de Taylor:
\begin{eqnarray}
k_2 &=&f(t^{(n)}+c_2h, \tilde{u}_2)\\
   &=&f(t^{(n)}+c_2h, u^{(n)}+ a_{21}h k_1)\\
   &=&f(t^{(n)}, u^{(n)}) + h\left[c_2 \frac{\partial f}{\partial t}+ a_{21} k_1\frac{\partial f}{\partial u}\right]+O(h^2)\\
   &=&k_1 + h\left(c_2 \frac{\partial f}{\partial t}+ a_{21} k_1\frac{\partial f}{\partial u}\right)+O(h^2)
   \end{eqnarray}
fazendo com que \eqref{eq:RK2a} se torne
\begin{eqnarray}
  u^{(n+1)}&=&u^{(n)}   + h b_1 k_1 + hb_2k_2 \\
         &=&u^{(n)}  + h b_1 k_1 + hb_2\left[ k_1 + h\left(c_2 \frac{\partial f}{\partial t}+ a_{21} k_1\frac{\partial f}{\partial u}\right)+O(h^2)\right]\\
         &=&u^{(n)}  + h (b_1+b_2) k_1 + h^2 b_2\left(c_2 \frac{\partial f}{\partial t}+ a_{21} k_1\frac{\partial f}{\partial u}\right)+O(h^3)\label{eq:rk2_1}
\end{eqnarray}
Usando a equação diferencial ordinária que desejamos resolver e derivando-a em $t$, obtemos:
\begin{eqnarray}
  u'(t)    &=&f(t,u(t)),\\
  u''(t) &=&\frac{\partial f}{\partial t}+\frac{\partial f}{\partial u} u'(t) = \frac{\partial f}{\partial t} +\frac{\partial f}{\partial u} f(t,u(t)).
\end{eqnarray}
Agora,  expandimos em série de Taylor a solução exata $u(t)$ em $t=t^{(n)}$,
\begin{eqnarray}
  u(t^{(n+1)})&=&u(t^{(n)}+h)=u^{(n)}  + hu'(t) +\frac{h^2}{2} u''(t) + O(h^3)\\
            &=&u^{(n)}  + hf(t,u^{(n)}) +\frac{h^2}{2}\frac{d }{d t}f(t,u(t))+O(h^3)\\
	    &=&u^{(n)}  + hf(t,u^{(n)}) +\frac{h^2}{2}\left(\frac{\partial f}{\partial t}+\frac{\partial f}{\partial u}u'(t^{(n)})\right) +O(h^3)\\
           &=&u^{(n)}  + hf(t,u^{(n)}) +\frac{h^2}{2}\left(\frac{\partial f}{\partial t}+\frac{\partial f}{\partial u}f(t,u^{(n)})\right) +O(h^3)\\
           &=&u^{(n)}  + hf(t,u^{(n)}) +\frac{h^2}{2}\left(\frac{\partial f}{\partial t}+\frac{\partial f}{\partial u}k_1\right) +O(h^3)\label{eq:rk2_2}
           \end{eqnarray}
Finalmente comparamos os termos em \eqref{eq:rk2_1} e \eqref{eq:rk2_2} de forma a haver concordância na expansão de Taylor até segunda ordem, isto é, restando apenas o erro de ordem 3 e produzindo um método de ordem de precisão $p=2$:
\begin{equation}\label{eq:cond_euler_2est}
  b_1+b_2=1, \quad b_2c_2 = \frac{1}{2} \quad \hbox{e} \quad a_{21}=c_2
\end{equation}
Este sistema é formado por três equações e quatro incógnitas, pelo que admite infinitas soluções. Para construir toda a família de soluções, escolha um parâmetro $\alpha\in (0,1] $ e defina a partir de \eqref{eq:cond_euler_2est}:
\begin{equation} b_1=1-\frac{1}{2\alpha},~~b_2=\frac{1}{2\alpha},~~c_2=\alpha,~~a_{21}=\alpha  \end{equation}
Portanto, obtemos o seguinte esquema genérico:
\begin{center}
\begin{tabular}{c|cc}
  \multirow{2}{*}{$c$}    &  \multicolumn{2}{c}{\multirow{2}{*}{$A$}} 	     \\
                          &  \multicolumn{2}{c}{}                            \\
  \hline
                          & \multicolumn{2}{c}{$~~b~~$}
\end{tabular}
~~~=~~~
\begin{tabular}{c|cc}
   &   &   \\
  $c_2$ & $a_{21}$ &   \\  \hline
      & $b_1$ &$b_2$
\end{tabular}
~~~=~~~
\begin{tabular}{c|cc}
   &   &   \\
  $\alpha$ & $\alpha$ &   \\  \hline
      & $\left(1-\frac{1}{2\alpha}\right)$ &$\frac{1}{2\alpha}$
\end{tabular},~~~$0<\alpha\leq 1$.
\end{center}
 Algumas escolhas comuns são $\alpha=\frac{1}{2}$, $\alpha=\frac{2}{3}$ e $\alpha=1$:
\begin{center}
\begin{tabular}{c|cc}
  $0$ &     &   \\
  $\frac{1}{2}$ & $\frac{1}{2}$ &   \\  \hline
      & $0$ & $1$
\end{tabular},\hphantom{~~~e~~~}
\begin{tabular}{c|cc}
  $0$ &   &   \\
  $\frac{2}{3}$ & $\frac{2}{3}$ &   \\  \hline
    & $\frac{1}{4}$ & $\frac{3}{4}$
\end{tabular} ~~~e~~~
\begin{tabular}{c|cc}
  $0$ &   &   \\
  $1$ & $1$ &   \\  \hline
      & $\frac{1}{2}$ &$\frac{1}{2}$
\end{tabular}.
\end{center}
Note que a tabela da direita fornece o método Euler modificado $\left(\alpha=1\right)$. O esquema iterativo assume a seguinte forma:
\begin{eqnarray}
  \tilde{u}_1 &=&u^{(n)} \\
  \tilde{u}_2 &=&u^{(n)}  +  h\alpha k_1 \\
  u^{(n+1)}&=&u^{(n)}  + h \left[ \left(1-\frac{1}{2\alpha}\right)k_1+\frac{1}{2\alpha} k_2\right],
\end{eqnarray}
onde $k_1=f(t^{(n)},u^{(n)})$ e $k_2=f(t^{(n)}+h\alpha,\tilde{u}_2)$. Ou, equivalentemente:
\begin{eqnarray}
k_1&=&f(t^{(n)},u^{(n)})\\
k_2&=&f(t^{(n)}+h\alpha,u^{(n)}  +  h\alpha k_1)\\
  u^{(n+1)}&=&u^{(n)}  + h \left[\left (1-\frac{1}{2\alpha}\right)k_1+\frac{1}{2\alpha} k_2\right],
\end{eqnarray}


\subsection{Métodos de Runge-Kutta com três estágios}\label{sec:RK_3_subsec}
Os métodos de Runge-Kutta com 3 estágios podem ser descritos na forma tabular como:
\begin{center}
\begin{tabular}{c|ccc}
  $0$              &          &               & \\
  $c_2$            & $a_{21}$ &               & \\
  $c_3$            & $a_{31}$ & $a_{32}$      & \\  \hline
                   & $b_{1}$  & $b_{2}$       & $b_{3}$
\end{tabular}
\end{center}


Seguindo um procedimento similar ao da Seção~\ref{eq:RK_2_subsec}, podemos obter as condições equivalentes às condições \eqref{eq:cond_euler_2est} para um método com $\nu =3$ e ordem $p=3$, as quais são:
\begin{eqnarray}\label{eq:cond_euler_3est}
  b_1+b_2+b_3		&=& 1,            		\label{eq:cond_euler_3est:1} \\
  b_2c_2+b_3c_3 	&=& \frac{1}{2},	 	\label{eq:cond_euler_3est:2} \\
  b_2c_2^2+b_3c_3^2	&=& \frac{1}{3}, 		\label{eq:cond_euler_3est:3} \\
  b_3a_{32}c_2		&=& \frac{1}{6},		\label{eq:cond_euler_3est:4} \\
  a_{21}		&=& c_2,			\label{eq:cond_euler_3est:5} \\
  a_{31}+a_{32}		&=& c_3.			\label{eq:cond_euler_3est:6}
\end{eqnarray}
Assim, temos 6 condições para determinar 8 incógnitas, o que implica a existência de uma enorme família de métodos de Runge-Kutta com três estágios e ordem $p=3$. Quando $c_2\neq c_3$ e os denominadores abaixo não se anulam, podemos fixar $c_2$ e $c_3$ e determinar os demais coeficientes:
\begin{eqnarray}
 b_1&=&1-\frac{1}{2c_2}-\frac{1}{2c_3}+\frac{1}{3c_2c_3}\\
 b_2&=&\frac{3c_3-2}{6c_2(c_3-c_2)}\\
 b_3&=&\frac{2-3c_2}{6c_3(c_3-c_2)}\\
 a_{21}&=&c_2\\
 a_{31}&=&c_3-\frac{1}{6b_3c_2}\\
 a_{32}&=&\frac{1}{6b_3c_2}
\end{eqnarray}


Um exemplo obtido pela família acima é o método clássico de Runge-Kutta de três estágios, com $c_2=\frac{1}{2}$ e $c_3=1$. O método de Nyström mostrado a seguir também tem três estágios e ordem três, mas corresponde ao caso especial $c_2=c_3=\frac{2}{3}$ e, por isso, não é obtido substituindo diretamente esses valores nas fórmulas que contêm o fator $c_3-c_2$:
\begin{center}
\begin{tabular}{c|ccc}
  $0$           &               &               & \\
  $\frac{1}{2}$ & $\frac{1}{2}$ &               & \\
  $1$           & $-1$          & $2$           & \\  \hline
                & $\frac{1}{6}$ & $\frac{4}{6}$ & $\frac{1}{6}$
\end{tabular}
~~~e~~~~~~\begin{tabular}{c|ccc}
  $0$ &     &   & \\
  $\frac{2}{3}$ & $\frac{2}{3}$ &   & \\
  $\frac{2}{3}$ & $0$ &$\frac{2}{3}$& \\  \hline
      & $\frac{2}{8}$ &$\frac{3}{8}$& $\frac{3}{8}$
\end{tabular}.
\end{center}

\subsection{Métodos de Runge-Kutta com quatro estágios}\label{ssec:RK_4_subsec}
As técnicas utilizadas nas seções \ref{eq:RK_2_subsec}  e \ref{sec:RK_3_subsec} podem ser usadas para obter métodos de quarta ordem ($p=4$) e quatro estágios ($\nu=4$). As seguintes tabelas descrevem os dois esquemas mais conhecidos de Runge-Kutta quarta ordem com quatro estágios. O primeiro é denominado \emph{método de Runge-Kutta 3/8} e o segundo é chamado de \emph{método de Runge-Kutta clássico}.

\begin{center}
\begin{tabular}{c|cccc}
  $0$ 		&     			&   		&   		&    \\
  $\frac{1}{3}$ & $\frac{1}{3}$ 	&   		&   		&    \\
  $\frac{2}{3}$ & $-\frac{1}{3}$ 	&$1$		&   		&    \\
  $1$ 		& $1$ 			&$-1$		&$1$		&    \\  \hline
		& $\frac{1}{8}$ 	&$\frac{3}{8}$	& $\frac{3}{8}$	& $\frac{1}{8}$
\end{tabular}
~~~\hbox{e}~~~~~~~
\begin{tabular}{c|cccc}
  $0$ 		&     			&   		&   		&    \\
  $\frac{1}{2}$ & $\frac{1}{2}$ 	&   		&   		&    \\
  $\frac{1}{2}$ & $0$ 			&$\frac{1}{2}$	&   		&    \\
  $1$ 		& $0$ 			&$0$		&$1$		&    \\  \hline
		& $\frac{1}{6}$ 	&$\frac{2}{6}$	& $\frac{2}{6}$	& $\frac{1}{6}$
\end{tabular}
\end{center}

O método de Runge-Kutta clássico é certamente o mais notório dos métodos de Runge-Kutta e seu esquema iterativo pode ser escrito como a seguir:
\begin{eqnarray}
k_1&=&f\left(t^{(n)},u^{(n)}\right)\\
k_2&=&f\left(t^{(n)}+h/2,u^{(n)}+hk_1/2\right)\\
k_3&=&f\left(t^{(n)}+h/2,u^{(n)}+hk_2/2\right)\\
k_4&=&f\left(t^{(n)}+h,u^{(n)}+hk_3\right)\\
u^{(n+1)}&=&u^{(n)}+\frac{h}{6}\left(k_1+2k_2+2k_3+k_4\right)
\end{eqnarray}
A seguinte heurística, usando o método de Simpson para quadratura numérica, pode ajudar a compreender os estranhos coeficientes:
\begin{eqnarray}
u({t^{(n+1)}})-u({t^{(n)}})&=&\int_{t^{(n)}}^{t^{(n+1)}}f(t,u(t))dt \\
&\approx& \frac{h}{6}\left[ f\left(t^{(n)},u(t^{(n)})\right)+4f\left(t^{(n)}+h/2,u(t^{(n)}+h/2)\right)\right.\\
&+&\left.f\left(t^{(n)}+h,u(t^{(n)}+h)\right)\right]\\
&\approx& \frac{h}{6}\left[k_1+4\left(\frac{k_2+k_3}{2}\right)+k_4\right]
\end{eqnarray}
onde $k_1$ e $k_4$ representam os valores de $f(t,u)$ nos extremos; $k_2$ e $k_3$ são duas aproximações diferentes para a inclinação no meio do intervalo.

\subsection*{Exercícios resolvidos}
\begin{exeresol}\label{exeresol:problema_resolvido_RK3} Construa o esquema iterativo o método clássico de Runge-Kutta três estágios cuja tabela é dada a seguir:

\begin{tabular}{c|ccc}
  $0$           &               &               & \\
  $\frac{1}{2}$ & $\frac{1}{2}$ &               & \\
  $1$           & $-1$          & $2$           & \\  \hline
                & $\frac{1}{6}$ & $\frac{4}{6}$ & $\frac{1}{6}$
\end{tabular}

\end{exeresol}
\begin{resol}
 \begin{eqnarray}
    k_1&=&f\left(t^{(n)},u^{(n)}\right)\\
    k_2&=&f\left(t^{(n)}+h/2,u^{(n)}+hk_1/2\right)\\
    k_3&=&f\left(t^{(n)}+h,u^{(n)}+h(-k_1+2k_2)\right)\\
  u^{(n+1)}&=&u^{(n)}+h\frac{k_1+4k_2+k_3}{6}
 \end{eqnarray}
\end{resol}

\begin{exeresol} Utilize o método clássico de Runge-Kutta três estágios para calcular o valor de $u(2)$ com passos $h=10^{-1}$ e $h=10^{-2}$ para o seguinte problema de valor inicial:
\begin{eqnarray}
 u'(t)&=& -u(t)^2 + t,\\
 u(0) &=&0.
\end{eqnarray}
\ifisscilab
Aplicando o processo iterativo obtido no Problema Resolvido~\ref{exeresol:problema_resolvido_RK3}, obtemos a seguinte rotina:
\begin{verbatim}
function [y]= f(t,u)
    y= t-u**2
endfunction


function [u] = RK3_classico(h,Tmax,u1)
    itmax = Tmax/h;
    u=zeros(itmax+1)
    u(1)=u1

    for i = 1:itmax
        t=(i-1)*h
        k1 = f(t, u(i))
        k2 = f(t+h/2, u(i) + h*k1/2)
        k3 = f(t+h,   u(i) + h*(2*k2-k1))

        u(i+1) = u(i) + h*(k1+4*k2+k3)/6
    end
endfunction
 \end{verbatim}
A qual pode ser invocada com:
\begin{verbatim}
->sol=RK3_classico(1e-2,2,0);sol(201)
 ans  =

    1.1935760016451

-->sol=RK3_classico(1e-3,2,0);sol(2001)
 ans  =

    1.1935759753635
\end{verbatim}


\fi
\ifisoctave
Aplicando o processo iterativo obtido no Problema Resolvido~\ref{exeresol:problema_resolvido_RK3}, obtemos a seguinte rotina:
\begin{verbatim}
function [y]= f(t,u)
  y= t-u**2;
endfunction

function [u] = RK3_classico(h,Tmax,u1)
  itmax = Tmax/h;
  u=zeros(itmax+1);
  u(1)=u1;

  for i = 1:itmax
    t=(i-1)*h;
    k1 = f(t, u(i));
    k2 = f(t+h/2, u(i) + h*k1/2);
    k3 = f(t+h,   u(i) + h*(2*k2-k1));

    u(i+1) = u(i) + h*(k1+4*k2+k3)/6;
  end
endfunction
 \end{verbatim}
A qual pode ser invocada como:
\begin{verbatim}
>> sol=RK3_classico(1e-2,2,0);sol(201)
ans =  1.19357600164514
>> sol=RK3_classico(1e-3,2,0);sol(2001)
ans =  1.19357597536351
\end{verbatim}
\fi

\ifispython
Aplicando o processo iterativo obtido no Problema Resolvido~\ref{exeresol:problema_resolvido_RK3}, obtemos a seguinte rotina:
\begin{verbatim}
 def f(t,u):
	return t-u**2


def RK3_classico(h,Tmax,u1):
  	itmax = int(round(Tmax/h))
	u=np.empty(itmax+1)
	u[0]=u1

	for i in range(itmax):
		t=i*h

		k1 = f(t,     u[i])
		k2 = f(t+h/2, u[i] + h*k1/2)
		k3 = f(t+h,   u[i] + h*(2*k2-k1))

		u[i+1] = u[i] + h*(k1+4*k2+k3)/6
	return u



Tmax=2			#tempo maximo de simulacao
u1=0			#condicoes iniciais na forma vetorial
h=1e-2			#passo

sol=RK3_classico(h,Tmax,u1);
itmax = int(round(Tmax/h))print(sol[itmax])
\end{verbatim}

\fi

\end{exeresol}

\subsection*{Exercícios}
\begin{exer} Aplique o esquema de Runge-Kutta segunda ordem com dois estágios cujos coeficientes são dados na tabela a seguir
 \begin{tabular}{c|cc}
  $0$ &   &   \\
  $\frac{2}{3}$ & $\frac{2}{3}$ &   \\  \hline
    & $\frac{1}{4}$ & $\frac{3}{4}$
\end{tabular}
para resolver o problema de valor inicial dado por:
\begin{eqnarray}
x'(t)&=&\sin(x(t)),\\
x(0)&=&2.
\end{eqnarray}
para $t=2$ com $h=1e-1$, $h=1e-2$ e $h=1e-3$. Expresse sua resposta com oito dígitos significativos corretos.
\end{exer}
\begin{resp}
 $2,9677921$, $2,9682284$ e $2,9682325$.
\end{resp}


\begin{exer}Resolva pelo método de Euler, Euler melhorado, Runge-Kutta clássico três estágios e Runge-Kutta clássico quatro estágios o problema de valor inicial tratados nos  exercícios resolvidos \ref{exeresol:exeresol1} e \ref{exeresol:exeresol1_euler_melhorado} dado por:
\begin{eqnarray}
     u'(t)&=& -0,5u(t)+2+t\\
            u(0) &=&  8
\end{eqnarray}

Usando os seguintes passos: $h=1$, $h=10^{-1}$, $h=10^{-2}$ e $h=10^{-3}$ e compare a solução aproximada em $t=1$ com as soluções obtidas com a solução exata dada por:
\begin{equation}
     u(t) = 2t+8e^{-t/2} \Longrightarrow u(1)=2+8e^{-1/2} \approx 6,85224527770107
\end{equation}
\end{exer}

\begin{resp}
\begin{center}
 \begin{tabular}{|l|l|l|l|l|}%\label{pvi:tab_euler}
\hline
Euler & 6,0000000 & 6,7898955 & 6,8461635 & 6,8516386\\
\hline
$\varepsilon_{rel}$ & 1,2e-01 & 9,1e-03 & 8,9e-04 & 8,9e-05\\
\hline
Euler mod, & 7,0000000 & 6,8532949 & 6,8522554 & 6,8522454\\
\hline
$\varepsilon_{rel}$ & 2,2e-02 & 1,5e-04 & 1,5e-06 & 1,5e-08\\
\hline
RK$_3$ & 6,8333333 & 6,8522321 & 6,8522453 & 6,8522453\\
\hline
$\varepsilon_{rel}$ & 2,8e-03 & 1,9e-06 & 1,9e-09 & 1,8e-12\\
\hline
RK$_4$ & 6,8541667 & 6,8522454 & 6,8522453 & 6,8522453\\
\hline
$\varepsilon_{rel}$ & 2,8e-04 & 1,9e-08 & 1,9e-12 & 1,3e-15\\
\hline
\end{tabular}
\end{center}

% Veja o gráfico da solução para $h=1, 0.5, 0.1, 0.05$:
% \begin{figure}
% \includegraphics[width=\textwidth]{euler.eps}
% \end{figure}
\end{resp}



\begin{exer}Aplique o método de Euler, o método de Euler melhorado, o método clássico de Runge-Kutta três estágios e o método clássico de Runge-Kutta quatro estágios para resolver o problema de valor inicial dado por
\begin{eqnarray}
u'&=& u + t\\
u(0)&=&1
\end{eqnarray}
com passo $h=1$, $h=10^{-1}$, $h=10^{-2}$ e $h=10^{-3}$  para obter aproximações para $u(1)$. Compare com a solução exata dada do problema dada por $u(t) =  2e^t-t-1$ através do erro relativo e observe a ordem de precisão do método. Expresse a sua resposta com oito dígitos significativos para a solução e 2 dígitos significativos para o erro relativo.
\end{exer}
\begin{resp}

\begin{center}
 \begin{tabular}{|l|l|l|l|l|}%\label{pvi:tab_euler}
\hline
Euler & 2,0000000 & 3,1874849 & 3,4096277 & 3,4338479\\
\hline
$\varepsilon_{rel}$ & 4,2e-01 & 7,2e-02 & 7,8e-03 & 7,9e-04\\
\hline
Euler mod & 3,0000000 & 3,4281617 & 3,4364737 & 3,4365628\\
\hline
$\varepsilon_{rel}$ & 1,3e-01 & 2,4e-03 & 2,6e-05 & 2,6e-07\\
\hline
RK$_3$ & 3,3333333 & 3,4363545 & 3,4365634 & 3,4365637\\
\hline
$\varepsilon_{rel}$ & 3,0e-02 & 6,1e-05 & 6,5e-08 & 6,6e-11\\
\hline
RK$_4$ & 3,4166667 & 3,4365595 & 3,4365637 & 3,4365637\\
\hline
$\varepsilon_{rel}$ & 5,8e-03 & 1,2e-06 & 1,3e-10 & 1,2e-14\\
\hline
\end{tabular}
\end{center}
\end{resp}



\section{Métodos de Runge-Kutta implícitos}\label{sec:sec_IRK}\index{método!de Runge-Kutta implícito}
Nos métodos de Runge-Kutta explícitos, cada estágio depende apenas de estágios já calculados. Em um método implícito, essa restrição é removida e o valor desconhecido do próprio estágio pode aparecer na equação que o define. Como consequência, a realização de um passo pode exigir a solução de uma equação algébrica, linear ou não linear.

O exemplo mais simples é a relação
\begin{equation}\label{eq:IRK:exemplo1}
  \begin{split}
    y^{(n+1)}&=y^{(n)}+h y^{(n+1)},\\
    y^{(1)}&=1,
  \end{split}
\end{equation}
que corresponde ao método de Euler implícito aplicado ao problema
\begin{eqnarray}
 y'(t)&=&y(t),\\
 y(0)&=&1.
\end{eqnarray}
A Equação~\eqref{eq:IRK:exemplo1} é implícita porque $y^{(n+1)}$ aparece dos dois lados. Neste caso, podemos isolá-lo diretamente.

\subsection{Método de Euler implícito}\index{método!de Euler implícito}\label{sec:euler_imp}
Considere o problema
\begin{eqnarray}
  u'(t)  &=& f(t,u(t)),\\
  u(t^{(1)}) &=& a.
\end{eqnarray}
Integrando de $t^{(n)}$ a $t^{(n+1)}$, obtemos
\begin{eqnarray}
u(t^{(n+1)})&=&u(t^{(n)})+\int_{t^{(n)}}^{t^{(n+1)}}f(t,u(t))\,dt.
\end{eqnarray}
No método de Euler implícito, aproximamos o integrando por seu valor no extremo direito do intervalo:
\begin{eqnarray}
u^{(n+1)}&=&u^{(n)}+h f(t^{(n+1)},u^{(n+1)}).
\end{eqnarray}

Se $f$ for linear em $u$, frequentemente podemos isolar $u^{(n+1)}$. Caso contrário, em geral é necessário resolver uma equação não linear em cada passo, por exemplo pelo método de Newton.

\begin{ex}
Considere o problema
\begin{eqnarray}
u'(t)&=&\lambda u(t),\\
u(0)&=&1.
\end{eqnarray}
O método de Euler implícito fornece
\begin{eqnarray}
 y^{(n+1)}&=&y^{(n)}+h\lambda y^{(n+1)},\\
 y^{(1)}&=&1.
\end{eqnarray}
Logo,
\begin{eqnarray}
 y^{(n+1)}&=&\frac{1}{1-\lambda h}y^{(n)},\\
 y^{(1)}&=&1.
\end{eqnarray}
\end{ex}

\subsection{O método trapezoidal}\index{método!trapezoidal}\label{sec:trapezoidal}
Aplicando a regra do trapézio à identidade integral,
\begin{eqnarray}
u(t^{(n+1)})&=&u(t^{(n)})+\int_{t^{(n)}}^{t^{(n+1)}}f(t,u(t))\,dt,
\end{eqnarray}
obtemos o \emph{método trapezoidal}
\begin{eqnarray}
u^{(n+1)}&=&u^{(n)}+\frac{h}{2}\left[f(t^{(n)},u^{(n)})+
f(t^{(n+1)},u^{(n+1)})\right],\\
u^{(1)}&=&a.
\end{eqnarray}
O método é implícito e tem ordem dois.

\begin{ex}
Para o problema
\begin{eqnarray}
u'(t)&=&\lambda u(t),\\
u(0)&=&1,
\end{eqnarray}
a recorrência é
\begin{eqnarray}
 y^{(n+1)}&=&y^{(n)}+\frac{\lambda h}{2}\left[y^{(n)}+y^{(n+1)}\right].
\end{eqnarray}
Isolando $y^{(n+1)}$,
\begin{eqnarray}
 y^{(n+1)}&=&\frac{1+\lambda h/2}{1-\lambda h/2}y^{(n)}.
\end{eqnarray}
\end{ex}

\subsection{O método theta}
Os métodos de Euler explícito, Euler implícito e trapezoidal pertencem à família de métodos theta:
\begin{eqnarray}
u^{(n+1)}&=&u^{(n)}+h\left[(1-\theta)f(t^{(n)},u^{(n)})
+\theta f(t^{(n+1)},u^{(n+1)})\right],
\end{eqnarray}
onde $0\leq\theta\leq1$. Para $\theta=0$, obtemos o método de Euler explícito; para $\theta=1/2$, o método trapezoidal; e para $\theta=1$, o método de Euler implícito.

\subsection*{Exercícios resolvidos}

\begin{exeresol}
Considere o problema de valor inicial
\begin{eqnarray}
y'(t)&=&y(t)(1-y(t)),\\
y(0)&=&\frac{1}{2}.
\end{eqnarray}
Construa a recursão via método de Euler implícito e explicite $y^{(n+1)}$.
\end{exeresol}
\begin{resol}
O método de Euler implícito produz
\begin{equation}
y^{(n+1)}=y^{(n)}+h y^{(n+1)}\left(1-y^{(n+1)}\right),
\end{equation}
ou, equivalentemente,
\begin{equation}
h\left[y^{(n+1)}\right]^2+(1-h)y^{(n+1)}-y^{(n)}=0.
\end{equation}
Pela fórmula da equação quadrática,
\begin{equation}
y^{(n+1)}=
\frac{-(1-h)\pm\sqrt{(1-h)^2+4hy^{(n)}}}{2h}.
\end{equation}
Como $y^{(n)}>0$, escolhemos a raiz positiva:
\begin{equation}
y^{(n+1)}=
\frac{-(1-h)+\sqrt{(1-h)^2+4hy^{(n)}}}{2h}.
\end{equation}
Para evitar cancelamento quando $h$ é pequeno, é preferível racionalizar a expressão:
\begin{equation}
y^{(n+1)}=
\frac{2y^{(n)}}{(1-h)+\sqrt{(1-h)^2+4hy^{(n)}}}.
\end{equation}
\end{resol}

\subsection*{Exercícios}

\construirExer

\section{O método de Taylor}
Uma maneira alternativa de aumentar a ordem dos métodos de Euler anteriormente descritos consiste em truncar a série de Taylor de $u(t+h)$:
\begin{eqnarray}
 u(t+h)=u(t) +h u'(t)+ \frac{h^2}{2!}u''(t)+\frac{h^3}{3!}u'''(t)+\ldots
\end{eqnarray}

Utilizando dois termos temos o método de Euler. Utilizando os três primeiros termos da série e substituindo $u'(t)=f(t,u)$ e $u''(t)=\frac{d f}{d t}(t,x)$ temos o \emph{método de Taylor de ordem $2$}
\begin{eqnarray}
   u^{(n+1)}&=&u^{(n)} +h f(t^{(n)},u^{(n)})\\&+& \frac{h^2}{2!} \left[ \frac{\partial}{\partial t}f(t^{(n)},u^{(n)})+\frac{\partial}{\partial u}f(t^{(n)},u^{(n)})f(t^{(n)},u^{(n)})\right]
\end{eqnarray}
onde usamos a regra da cadeia para obter:
\begin{eqnarray}
  \frac{d }{d t}f(t,u)= \frac{\partial}{\partial t}f(t,u)+\frac{\partial}{\partial u}f(t,u)u'(t)=\frac{\partial}{\partial t}f(t,u)+\frac{\partial}{\partial u}f(t,u)f(t,u)
\end{eqnarray}



De forma análoga, o método de Taylor de ordem $3$ é
\begin{eqnarray}
   u^{(n+1)}=u^{(n)} +h f(t^{(n)},u^{(n)})+ \frac{h^2}{2!}\frac{d f}{d t}(t^{(n)},u^{(n)})+\frac{h^3}{3!}\frac{d^2 f}{d t^2}(t^{(n)},u^{(n)})
\end{eqnarray}

\subsection*{Exercícios resolvidos}

\construirExeresol

\subsection*{Exercícios}

\construirExer

%\section{Métodos de passo múltiplo}
\section{Método de Adams-Bashforth}\index{Método!de Adams-Bashforth}\index{Método!de passo múltiplo}\label{sec:Adams_Bashforth}
Seja o problema de valor inicial
\begin{eqnarray}
  u'(t) &= f(t,u(t)) \\
  u(t_0) &= a
\end{eqnarray}

Nos métodos de passo simples, os valores calculados para $f\left(t^{(n)},u(t^{(n)})\right)$ nos passos anteriores são desprezados ao calcular o próximo passo. Nos métodos de passo múltiplo, os valores de $f\left(t,u)\right)$, nos passos $n$, $n+1$, ..., $n+s-1$ são utilizados ao calcular $f$ em $t^{(n+s)}$.

Integrando a equação diferencial no intervalo $[t^{(n+s-1)},t^{(n+s)}]$, obtemos:
\begin{equation}\label{eq:mpm1}
  u^{(n+s)}  = u^{(n+s-1)}  + \int_{t^{(n+s-1)}}^{t^{(n+s)}} f(t,u(t)) dt
\end{equation}
No método de Adams-Bashforth, o integrando em \eqref{eq:mpm1} é aproximado pelo polinômio que interpola $f(t^{(k)},u^{(k)})$ para $k = n, n+1, n+2, \ldots, n+s-1$, isto é:
\begin{eqnarray}%\label{pvi:mpm1}
  u^{(n+s)}  &=& u^{(n+s-1)}  + \int_{t^{(n+s-1)}}^{t^{(n+s)}} p(t) dt
\end{eqnarray}
onde $p(t)$ é polinômio de grau $s-1$ dado na forma de Lagrange por:
\begin{equation} p(t)=\sum_{j=0}^{s-1}\left[f(t^{(n+j)},u^{(n+j)}) \prod_{k=0,k\neq j}^{s-1} \frac{t-t^{(n+k)}}{t^{(n+j)}-t^{(n+k)}}\right] \end{equation}
Agora observamos que
\begin{equation} \int _{t^{(n+s-1)}}^{t^{(n+s)}} p(t) dt=h\sum_{j=0}^{s-1}\beta_j f(t^{(n+j)},u^{(n+j)}) \end{equation}
onde
\begin{equation}\label{eq:betaj}
\beta_j= \frac{1}{h}\int_{t^{(n+s-1)}}^{t^{(n+s)}} \prod_{k=0,k\neq j}^{s-1} \frac{t-t^{(n+k)}}{t^{(n+j)}-t^{(n+k)}}dt
\end{equation}

e obtemos a relação de recorrência:
\begin{eqnarray}\label{eq:mpm2}
  u^{(n+s)}  &=& u^{(n+s-1)}  + h\sum_{j=0}^{s-1}\beta_j f(t^{(n+j)},u^{(n+j)})
\end{eqnarray}
Observe que a integral envolvida no cálculo dos coeficientes $\beta_j$ em \eqref{eq:betaj} pode ser simplificada via a mudança de variáveis $t=t^{(n+s-1)}+h\tau$:
\begin{eqnarray}\label{eq:betaj_s}
\beta_j&=& \int_{0}^{1} \prod_{k=0,k\neq j}^{s-1} \frac{\tau+s-k-1}{j-k}d\tau\\
&=& \frac{(-1)^{s-j-1}}{j!(s-j-1)!}\int_{0}^{1} \prod_{k=0,k\neq j}^{s-1}(\tau+s-k-1)d\tau\\
&=& \frac{(-1)^{s-j-1}}{j!(s-j-1)!}\int_{0}^{1} \prod_{k=0,k\neq s-j-1}^{s-1}(\tau+k)d\tau
\end{eqnarray}


\begin{obs}[Ordem do método de Adams-Bashforth] Da teoria de interpolação (ver capítulos \ref{cap:interp} e \ref{cap:integracao}), temos que o erro de aproximação da integral de uma função suficientemente suave por um polinômio interpolador em $s$ pontos é de ordem $s+1$. Assim, o erro local de truncamento do método de Adams-Bashforth com $s$ passos é $s+1$ e, portanto, o erro global de truncamento é de ordem $s$.
\end{obs}


\begin{ex} Calcule os coeficientes de Adams-Bashforth para $s=2$ e, depois, construa seu processo iterativo.
 \begin{eqnarray}
  \beta_0&=&-\int_0^1 (\tau+2-1-1)d\tau=-\frac{1}{2}\\
  \beta_1&=&\int_0^1 (\tau+2-0-1)d\tau=\frac{3}{2}
 \end{eqnarray}
 O processo iterativo é dado por:
 \begin{equation}
  y^{(n+2)}=y^{(n+1)}+\frac{h}{2}\left[3f\left(t^{(n+1)},y^{(n+1)}\right)-f\left(t^{(n)},y^{(n)}\right)\right]
 \end{equation}
\end{ex}


\begin{ex} Calcule os coeficientes e de Adams-Bashforth para $s=3$ e, depois, construa o processo iterativo.
 \begin{eqnarray}
  \beta_0&=&\frac{1}{2}\int_0^1 {(\tau+3-1-1)}\cdot {(\tau+3-2-1)} d\tau=\frac{5}{12}\\
  \beta_1&=&-\int_0^1 {(\tau+3-0-1)}\cdot {(\tau+3-2-1)} d\tau=-\frac{4}{3}\\
  \beta_2&=&\frac{1}{2}\int_0^1 {(\tau+3-0-1)}\cdot {(\tau+3-1-1)} d\tau=\frac{23}{12}
 \end{eqnarray}
 O processo iterativo é dado por:
 \begin{equation}
  y^{(n+3)}=y^{(n+2)}+\frac{h}{12}\left[23f\left(t^{(n+2)},y^{(n+2)}\right)-16f\left(t^{(n+1)},y^{(n+1)}\right)+5f\left(t^{(n)},y^{(n)}\right)\right]
 \end{equation}
\end{ex}

\begin{obs} Os coeficientes do método de Adams-Bashforth de ordem $s$ podem, alternativamente, ser obtidos exigindo que o sistema seja exato para $f(t,u)=t^0$, $f(t,u)=t^1$, $f(t,u)=t^2$, \ldots, $f(t,u)=t^{s-1}$.

\end{obs}

\begin{ex} Obtenha o método de Adams-Bashforth para $s=4$ como
\begin{eqnarray}
  u^{(n+4)}  &=& u^{(n+3)}  + \int _{t^{(n+3)}}^{t^{(n+4)}} f(t,u(t)) dt \\
  u^{(n+4)}  &=& u^{(n+3)}  + h \sum_{m=0}^{3} b_m f^{(n+m)} \\
  u^{(n+4)}  &=& u^{(n+3)}  + h \left[b_3f^{(n+3)} +b_2f^{(n+2)} +b_1f^{(n+1)} +b_0f^{(n)}\right]
\end{eqnarray}
Para isso devemos obter $[b_3,b_2,b_1,b_0]$ tal que o método seja exato para polinômios até ordem $3$. Podemos obter esses coeficientes de maneira análoga a obter os coeficientes de um método para integração.

Supondo que os nós $t^{(k)}$ estejam igualmente espaçados, e para facilidade dos cálculos, como o intervalo de integração é $[t^{(n+3)},t^{(n+4)}]$, translade $t^{(n+3)}$ para a origem tal que $[t^{(n)},t^{(n+1)},\ldots ,t^{(n+4)}]=[-3h,-2h,-h,0,h]$.

Considere a base $[1, t, t^2, t^3]$ e substitua $f(t)$ por cada um dos elementos desta basa, obtendo:
\begin{eqnarray}
  \int _0^{h} 1  \;dt = h             &=& h\left[ b_0(1)  +b_1(1)    + b_2(1)   + b_3(1)    \right]\\
  \int _0^{h} t  \;dt = \frac{h^2}{2}  &=& h\left[ b_0(-3h)  +b_1(-2h)   + b_2(-h) + b_3(0)  \right]\\
  \int _0^{h} t^2 \;dt = \frac{h^3}{3}  &=& h\left[b_0(-3h)^2 +b_1(-2h)^2  + b_2(-h)^2+ b_3(0)^2 \right]\\
  \int _0^{h} t^3 \;dt = \frac{h^4}{4} &=& h\left[ b_0(-3h)^3 +b_1(-2h)^3  + b_2(-h)^3+ b_3(0)^3 \right]
\end{eqnarray}
que pode ser escrito na forma matricial
\begin{eqnarray}
\left(
  \begin{array}{cccc}
    1  &  1  &  1  &  1\\
    -3  &  -2  &  -1  &  0\\
    9  &  4  &  1  &  0\\
    -27  &  -8  & -1  &  0
  \end{array}
\right)
\left(\begin{array}{c}  b_0 \\ b_1\\ b_2\\b_3   \end{array}\right)
=
\left(\begin{array}{c}  1  \\ 1/2 \\ 1/3 \\ 1/4  \end{array}\right)
\end{eqnarray}
Resolvendo o sistema obtemos
\begin{equation} [b_0,b_1,b_2,b_3]=\left[-\frac{9}{24},\frac{37}{24},-\frac{59}{24},\frac{55}{24}\right] \end{equation}
fornecendo o \emph{método de Adams-Bashforth de $4$ passos}
\begin{eqnarray}\label{AB4}
  u^{(n+4)}  &= u^{(n+3)}  + \frac{h}{24} [55 f^{(n+3)} -59f^{(n+2)} +37f^{(n+1)} -9f^{(n)}]
\end{eqnarray}
\end{ex}


A tabela abaixo mostra os coeficientes do método de Adams-Bashforth para até 8 passos.
 \begin{center}
\begin{tabular}{|c|cccccccc|}
\hline
1 & 1 &&&&&&&\\
2 & $-\frac{1}{2}$ & $\frac{3}{2}$&&&&&&\\
3 & $\frac{5}{12}$ & $-\frac{4}{3}$ & $\frac{23}{12}$&&&&&\\
4 & $-\frac{9}{24}$ & $\frac{37}{24}$ & -$\frac{59}{24}$ & $\frac{55}{24}$&&&&\\
5 & $\frac{251}{720}$ & -$\frac{637}{360}$ & $\frac{109}{30}$& -$\frac{1387}{360}$& $\frac{1901}{720}$&&&\\
6 & $-{\frac {95}{288}}$ & ${\frac {959}{480}}$ & $-{\frac {3649}{720}}$ & ${\frac {4991}{720}}$ & $-{\frac {2641}{480}}$ & $\frac {4277}{1440}$&&\\
7 & $\frac {19087}{60480}$ & $-{\frac {5603}{2520}}$ & ${\frac {135713}{20160}}$ & $-{\frac {10754}{945}}$ & ${\frac {235183}{20160}}$ & $-{\frac {18637}{2520}}$ & ${\frac {198721}{60480}}$&\\
8 & $-{\frac {5257}{17280}}$&${\frac {32863}{13440}}$&$-{\frac {115747}{13440}}$&${\frac {2102243}{120960}}$&$-{\frac {296053}{13440}}$&${\frac {242653}{13440}}$&$-{\frac {1152169}{120960}}$&${\frac {16083}{4480}}$\\
\hline
\end{tabular}
 \end{center}
\begin{obs} Note que os métodos de passo múltiplo requerem o conhecimento dos $s$ valores previamente computados para calcular $y^{(n+s)}$. Assim, para inicializar um algoritmo com mais de um passo, não é suficiente conhecer a condição inicial. Usualmente, calculam-se os valores iniciais necessários usando um algoritmo de passo simples da mesma ordem do método múltiplo passo a ser aplicado.
\end{obs}




\subsection*{Exercícios resolvidos}
\begin{exeresol}\label{exeresol:Adams_Bashforth:prob1} Resolva numericamente o problema de valor inicial dado por:
\begin{eqnarray}
 y'(t)&=& \sqrt{1+y(t)}\\
 y(0)&=& 0
\end{eqnarray}
aplicando o método de Adams-Bashforth de dois passos e inicializando o método através do método de Euler modificado. Calcule o valor de $y(1)$ com passo de tamanho $h=0,1$.
\end{exeresol}
\begin{resol}
Primeiro observamos que o processo recursivo do método de Adams é dado por:
\begin{equation}
  y^{(n+2)}=y^{(n+1)}+\frac{h}{2}\left[3f\left(t^{(n+1)},u(t^{(n+1)})\right)-f\left(t^{(n)},u(t^{(n)})\right)\right],~n=1,2,\ldots
 \end{equation}
O valor inicial é dado por $y^{(1)}=0$. No entanto, para inicializar o método, precisamos calcular $y^{(2)}$, para tal, aplicamos o método de Euler modificado:
\begin{eqnarray}
 k_1&=&\sqrt{1+0}=1\\
 k_2&=&\sqrt{1+0,1}=\sqrt{1,1}\approx 1,0488088 \\
 y^{(2)}&=& \frac{0,1}{2}\left(1+1,0488088 \right)=0,10244044
  \end{eqnarray}
Aplicando o método de Adams-Bashforth, obtemos:
  \begin{eqnarray}
 y^{(1)}&=&0\\
 y^{(2)}&=&0,10244044\\
 y^{(3)}&=&0,20993619\\
 y^{(4)}&=&0,32243326\\
 y^{(5)}&=&0,43993035\\
 y^{(6)}&=&0,56242745\\
 y^{(7)}&=&0,68992455\\
 y^{(8)}&=&0,82242165\\
 y^{(9)}&=&0,95991874\\
 y^{(10)}&=&1,10241584\\
 y^{(11)}&=&1,24991294
  \end{eqnarray}
\end{resol}

\ifispython
A seguinte rotina implementa o método:
\begin{verbatim}
def f(t,u):
	return np.sqrt(1+u)

def adams_bash_2(h,Tmax,u1):
	dim=np.size(u1)
	itmax=int(round(Tmax/h))
	u=np.empty((itmax+1,dim))
	u[0,:]=u1

	#inicaliza com RK2
	k1 = f(0,   u[0,:])
	k2 = f(h, u[0,:] + k1* h)
	u[1,:] = u[0,:] + (k1+k2)* h/2

	fn_0=k1
	for i in range(itmax-1):
		t=(i+1)*h
		fn_1 = f(t,   u[i+1,:])
		u[i+2,:] = u[i+1,:] + h*(-.5*fn_0 + 1.5*fn_1)
		fn_0=fn_1
	return u





u0=0
h=1e-1
Tmax=1
u=adams_bash_2(h,Tmax,u0)

print(u)
\end{verbatim}
\fi

\emconstrucao

\subsection*{Exercícios}

\emconstrucao

\section{Método de Adams-Moulton}\index{Método!de Adams-Moulton}\label{sec:Adams_Moulton}
O método de Adams-Moulton, assim como o método de Adams-Bashforth, é um método de passo múltiplo. A diferença entre estes dois métodos é que Adams-Bashforth é explícito, enquanto Adams-Moulton é implícito, isto é, os valores de $f\left(t,u)\right)$, nos passos $n$, $n+1$, ..., $n+s-1$ e, inclusive, $n+s$ são utilizados ao calcular $f$ em $t^{(n+s)}$.

Considere o problema de valor inicial
\begin{eqnarray}
  u'(t) &= f(t,u(t)) \\
  u(t_0) &= a
\end{eqnarray}

Integrando a equação diferencial no intervalo $[t^{(n+s-1)},t^{(n+s)}]$, obtemos:
\begin{equation}\label{eq:mpm2a}
  u^{(n+s)}  = u^{(n+s-1)}  + \int_{t^{(n+s-1)}}^{t^{(n+s)}} f(t,u(t)) dt
\end{equation}
 Agora o integrando em \eqref{eq:mpm2a} é aproximado pelo polinômio que interpola $f(t^{(k)},u^{(k)})$ para $k = n, n+1, n+2, \ldots, n+s$, isto é:
\begin{equation}\label{eq:mpm2b}
  u^{(n+s)} = u^{(n+s-1)}  + \int_{t^{(n+s-1)}}^{t^{(n+s)}} p(t) dt
\end{equation}
onde $p(t)$ é polinômio de grau $s$ dado na forma de Lagrange por:
\begin{equation} p(t)=\sum_{j=0}^{s}\left[f(t^{(n+j)},u^{(n+j)}) \prod_{k=0,k\neq j}^{s} \frac{t-t^{(n+k)}}{t^{(n+j)}-t^{(n+k)}}\right] \end{equation}
Agora observamos que
\begin{equation} \int _{t^{(n+s-1)}}^{t^{(n+s)}} p(t) dt=h\sum_{j=0}^{s}\beta_j f(t^{(n+j)},u^{(n+j)}) \end{equation}
onde
\begin{equation}
\beta_j= \frac{1}{h}\int_{t^{(n+s-1)}}^{t^{(n+s)}} \prod_{k=0,k\neq j}^{s} \frac{t-t^{(n+k)}}{t^{(n+j)}-t^{(n+k)}}dt
\end{equation}
Aplicando a mudança de variáveis $t=t^{(n+s-1)}+h\tau$, temos:
\begin{eqnarray}\label{eq:betaj_m}
\beta_j&=& \int_0^1 \prod_{k=0,k\neq j}^{s} \frac{\tau+s-k-1}{j-k}d\tau\\
&=& \frac{(-1)^{s-j}}{j!(s-j)!}\int_{0}^{1} \prod_{k=0,k\neq j}^{s}(\tau+s-k-1)d\tau\\
&=& \frac{(-1)^{s-j}}{j!(s-j)!}\int_{0}^{1} \prod_{k=0,k\neq s-j-1}^{s}(\tau+k-1)d\tau
\end{eqnarray}

Assim, obtemos a relação de recorrência:
\begin{eqnarray}\label{eq:mam}
  u^{(n+s)}  &=& u^{(n+s-1)}  + h\sum_{j=0}^{s}\beta_j f(t^{(n+j)},u^{(n+j)})
\end{eqnarray}


\begin{obs} Os coeficientes do método de Adams-Moulton de $s$ passos podem, alternativamente, ser obtidos exigindo que o sistema seja exato para $f(t,u)=t^0$, $f(t,u)=t^1$, $f(t,u)=t^2$, \ldots, $f(t,u)=t^{s}$.

\end{obs}

\begin{ex} Obtenha o método de Adams-Moulton para $s=3$ na forma
\begin{eqnarray}\label{AM4}
u^{(n+3)}&=&u^{(n+2)}+\int_{t^{(n+2)}}^{t^{(n+3)}}f(t,u(t))\,dt\\
&\approx&u^{(n+2)}+h\left[b_0f^{(n+3)}+b_1f^{(n+2)}
+b_2f^{(n+1)}+b_3f^{(n)}\right].
\end{eqnarray}
Transladando $t^{(n+2)}$ para a origem, os quatro nós usados na interpolação são
\begin{equation}
[-2h,-h,0,h],
\end{equation}
associados, respectivamente, a $f^{(n)}$, $f^{(n+1)}$, $f^{(n+2)}$ e $f^{(n+3)}$. Para obter os pesos, exigimos exatidão para $1,t,t^2,t^3$. Na ordem $[f^{(n+3)},f^{(n+2)},f^{(n+1)},f^{(n)}]$, o sistema é
\begin{eqnarray}
\left(
\begin{array}{cccc}
1&1&1&1\\
1&0&-1&-2\\
1&0&1&4\\
1&0&-1&-8
\end{array}
\right)
\left(
\begin{array}{c}
b_0\\b_1\\b_2\\b_3
\end{array}
\right)
=
\left(
\begin{array}{c}
1\\1/2\\1/3\\1/4
\end{array}
\right).
\end{eqnarray}
Resolvendo o sistema,
\begin{equation}
[b_0,b_1,b_2,b_3]
=
\left[\frac{9}{24},\frac{19}{24},-\frac{5}{24},\frac{1}{24}\right].
\end{equation}
Portanto,
\begin{eqnarray}
u^{(n+3)}
&=&u^{(n+2)}
+\frac{h}{24}\left[9f^{(n+3)}+19f^{(n+2)}-5f^{(n+1)}+f^{(n)}\right].
\end{eqnarray}
\end{ex}




A tabela abaixo mostra os coeficientes do método de Adams-Moulton para até oito pontos de interpolação. Em cada linha, os coeficientes estão ordenados do valor mais antigo para o valor futuro.
 \begin{center}
\begin{tabular}{|c|cccccccc|}
\hline
1 & 1 &&&&&&&\\
2 & $\frac{1}{2}$ & $\frac{1}{2}$&&&&&&\\
3 & $-\frac{1}{12}$&$\frac{2}{3}$&$\frac{5}{12}$&&&&&\\
4 & $\frac{1}{24}$ & $-{\frac {5}{24}}$ & ${\frac {19}{24}}$ & $\frac{3}{8}$ &&&&\\
5 & $-{\frac {19}{720}}$ & ${\frac {53}{360}}$ & $-{\frac {11}{30}}$ & ${\frac {323}{360}}$ & ${\frac {251}{720}}$&&&\\
6 & ${\frac {3}{160}}$ & $-{\frac {173}{1440}}$ & ${\frac {241}{720}}$ & $-{\frac {133}{240}}$ & ${\frac {1427}{1440}}$ & ${\frac {95}{288}}$&&\\
7 & $-{\frac {863}{60480}}$ & ${\frac {263}{2520}}$ & $-{\frac {6737}{20160}}$ & ${\frac {586}{945}}$ & $-{\frac {15487}{20160}}$ & ${\frac {2713}{2520}}$ &${\frac {19087}{60480}}$&\\
8 &$ {\frac {275}{24192}}$ & $-{\frac {11351}{120960}}$&${\frac {1537}{4480}}$&$-{\frac {88547}{120960}}$&${\frac {123133}{120960}}$&$-{\frac {4511}{4480}}$&${\frac {139849}{120960}}$&${\frac {5257}{17280}}$\\
\hline
\end{tabular}
 \end{center}

\begin{ex} O esquema iterativo de Adams-Moulton com três passos, isto é, $s=2$ é dado na forma:
 \begin{equation}
  u^{(n+2)}=u^{(n+1)}+\frac{h}{12}\left[5f\left(t^{(n+2)},u(t^{(n+2)})\right)+8f\left(t^{(n+1)},u(t^{(n+1)})\right)-f\left(t^{(n)},u(t^{(n)})\right)\right]
 \end{equation}
 \end{ex}


\subsection*{Exercícios resolvidos}
\begin{exeresol} Resolva o problema de valor inicial dado por:
\begin{eqnarray}
 u'(t)&=& -2u(t) + te^{-t}\\
 u(0)&=&-1
\end{eqnarray}
via Adams-Moulton com $s=2$ (três passos) com $h=0,1$ e $h=0,01$ e compare com a solução exata dada por $u(t)=(t-1)e^{-t}$ nos instantes $t=1$ e $t=2$. Inicialize com Euler modificado.

 \end{exeresol}
\begin{resol}
 Primeiro observamos que $f(t,u)=-2u + te^{-t}$ e que o esquema de Adams-Moulton pode ser escrito como:
 \begin{equation}
  u^{(n+2)}=u^{(n+1)}+\frac{h}{12}\left[5f\left(t^{(n+2)},u(t^{(n+2)})\right)+8f\left(t^{(n+1)},u(t^{(n+1)})\right)-f\left(t^{(n)},u(t^{(n)})\right)\right]
 \end{equation}
 de forma que:
 \begin{eqnarray}
  u^{(n+2)}&=&u^{(n+1)}+\frac{h}{12}\left[8f\left(t^{(n+1)},u(t^{(n+1)})\right)-f\left(t^{(n)},u(t^{(n)})\right)\right]+\frac{5h}{12}f\left(t^{(n+2)},u(t^{(n+2)})\right)\\
  &=&u^{(n+1)}+\frac{h}{12}\left[8f^{(n+1)}-f^{(n)}\right]+\frac{5h}{12}\left(t^{(n+2)}e^{-t^{(n+2)}} -2 u^{(n+2)}\right)
   \end{eqnarray}
 Assim:
 \begin{eqnarray}
  \left(1+\frac{5h}{6}\right)u^{(n+2)}
  &=&u^{(n+1)}+\frac{h}{12}\left[8f^{(n+1)}-f^{(n)}\right]+\frac{5h}{12}t^{(n+2)}e^{-t^{(n+2)}}
  \end{eqnarray}

 Os valores obtidos são:

 \begin{tabular}{|c|c|c|}
 \hline
  &t=1&t=2\\
 \hline
  h=0,1  & -0,000223212480142 & 0,135292280956\\
  h=0,01 & -2,02891229566e-07& 0,135335243537\\
  Exato  & 0 & 0,135335283237\\
  \hline
 \end{tabular}
\ifispython
A seguinte rotina implementa a recursão:
\begin{verbatim}
## resolve u'(t)=l*u(t) + g

def g(t):
	return t*np.exp(-t)




u0=-1
h=1e-2
Tmax=2
itmax=int(round(Tmax/h))

u=np.empty(itmax+1)
fn=np.empty(itmax+1)

u[0]=u0
l=-2
#Iniciliza com Euler modificado
k1= l*u[0] + g(0)
k2= l*(u[0]+h*k1) + g(h)
u[1]= u[0]+ h *(k1+k2)/2

fn[0]= k1
fn[1]= l*u[1] + g(h)



for n in range(itmax-1):
	gn2=g((n+2)*h)

	u[n+2]= (u[n+1] + h/12*(8*fn[n+1]-fn[n]) + 5*h/12*gn2 ) / (1+5*h/6)
	fn[n+2]=l*u[n+2]+gn2

for n in range(itmax+1):
	print(h*n, u[n], (h*n-1)*np.exp(-h*n))
\end{verbatim}

\fi
\ifisoctave
A seguinte rotina implementa a recursão:
\begin{verbatim}
function y = g(x)
  y=x*exp(-x);
endfunction

function u=moulton(u0,h,Tmax)
  itmax=Tmax/h;

  u=[0:h:Tmax];
  fn=[0:h:Tmax];

  u(1)=u0;
  l=-2;
  %Iniciliza com Euler modificado
  k1= l*u(1) + g(0);
  k2= l*(u(1)+h*k1) + g(h);
  u(2)= u(1)+ h *(k1+k2)/2;

  fn(1)= k1;
  fn(2)= l*u(2) + g(h);



  for n=1:itmax-1
    gn2=g((n+1)*h);

    u(n+2)= (u(n+1) + h/12*(8*fn(n+1)-fn(n)) + 5*h/12*gn2 ) / (1+5*h/6);
    fn(n+2)=l*u(n+2)+gn2;
  endfor

endfunction
\end{verbatim}

\fi
\ifisscilab
A seguinte rotina implementa a recursão:
\begin{verbatim}
function y = g(x)
  y=x*exp(-x);
endfunction

function u=moulton(u0,h,Tmax)
  itmax=Tmax/h;

  u=[0:h:Tmax];
  fn=[0:h:Tmax];

  u(1)=u0;
  l=-2;
  Iniciliza com Euler modificado
  k1= l*u(1) + g(0);
  k2= l*(u(1)+h*k1) + g(h);
  u(2)= u(1)+ h *(k1+k2)/2;

  fn(1)= k1;
  fn(2)= l*u(2) + g(h);



  for n=1:itmax-1
    gn2=g((n+1)*h);

    u(n+2)= (u(n+1) + h/12*(8*fn(n+1)-fn(n)) + 5*h/12*gn2 ) / (1+5*h/6);
    fn(n+2)=l*u(n+2)+gn2;
  end

endfunction

\end{verbatim}


\fi

\end{resol}




\begin{exeresol}\label{exeresol:Adams_Moulton:prob1} Repita o Problema~\ref{exeresol:Adams_Bashforth:prob1} pelo método de Adams-Moulton, isto, é  resolva numericamente o problema de valor inicial dado por:
\begin{eqnarray}
 y'(t)&=& \sqrt{1+y(t)},\\
 y(0)&=& 0,
\end{eqnarray}
aplicando o método de Adams-Moulton de dois passos. Calcule o valor de $y(1)$ com passo de tamanho $h=0,1$.
\end{exeresol}
\begin{resol}
Primeiro observamos que o processo recursivo do método de Adams é dado por:
\begin{equation}\label{eq:Adams_Moulton:prob1:rec}
  y^{(n+1)}=y^{(n)}+\frac{h}{2}\left[f\left(t^{(n+1)},u(t^{(n+1)})\right)+f\left(t^{(n)},u(t^{(n)})\right)\right],~n=1,2,\ldots
 \end{equation}
O valor inicial é dado por $y^{(1)}=0$. Primeiramente, precisamos isolar $y^{(n+1)}$ na Equação~\ref{eq:Adams_Moulton:prob1:rec}:
\begin{equation}
  y^{(n+1)}=y^{(n)}+\frac{h}{8} \sqrt {{h}^{2}+16+16y^{(n)}+8h\sqrt {1+y^{(n)}}} +\frac{h}{2}\sqrt {1+y^{(n)}}+\frac{h^2}{8},~n=1,2,\ldots
 \end{equation}

  \begin{eqnarray}
 y^{(1)}&=&0\\
 y^{(2)}&=&0,1025\\
y^{(3)}&=&0,21\\
y^{(4)}&=&0,3225\\
y^{(5)}&=&0,44\\
y^{(6)}&=&0,5625\\
y^{(7)}&=&0,69\\
y^{(8)}&=&0,8225\\
y^{(9)}&=&0,96\\
y^{(10)}&=&1,1025\\
y^{(11)}&=&1,25
\end{eqnarray}

\end{resol}


\begin{exeresol}\label{exeresol:exersol_moulton_nl} Resolva o problema de valor inicial dado por
\begin{eqnarray}
 y'(t)&=& y^3-y+t,\\
 y(0)&=& 0,
\end{eqnarray}
aplicando o método de Adams-Moulton de dois passos. Calcule o valor de $y(1)$ com passo de tamanho $h=0,1$ e $h=0,01$.
Primeiro observamos que o processo recursivo do método de Adams é dado por:
\begin{eqnarray}
  y^{(n+1)}&=&y^{(n)}+\frac{h}{2}\left[f\left(t^{(n+1)},u(t^{(n+1)})\right)+f\left(t^{(n)},u(t^{(n)})\right)\right],~n=1,2,\ldots,\\
  y^{(1)}&=&0
 \end{eqnarray}
Observamos que o problema de isolar $y^{(n+1)}$ pode ser escrito como
\begin{eqnarray}
  y^{(n+1)}-\frac{h}{2}f\left(t^{(n+1)},u(t^{(n+1)})\right) &=&y^{(n)}+\frac{h}{2}f\left(t^{(n)},u(t^{(n)})\right)
 \end{eqnarray}
O termo da esquerda é uma expressão não linear em $y^{(n+1)}$ e o termo da direita é conhecido, isto é, pode ser calculado com base nos valores anteriormente calculados. Devemos, então, escolher um método numérico de solução de equações algébricas não lineares (veja capítulo~\ref{cap:equacao1d}) como o método de Newton visto na Seção~\ref{sec:metodo_newton_1d} para resolver uma equação do tipo:
\begin{equation} u-\frac{h}{2}f(t^{(n+1)},u)=a, \end{equation}
isto é:
\begin{equation} u-\frac{h}{2}\left(u^3-u+t^{(n+1)}\right)=a, \end{equation}
com $a=y^{(n)}+\frac{h}{2}f\left(t^{(n)},u(t^{(n)})\right)$
Os valores obtidos são: $0,37496894$ e $0,37512382$ quando o método é inicializado com Euler melhorado.
\ifispython
A seguinte rotina implementa a recursão:
\begin{verbatim}
def f(t,u):
	return u**3-u+t

def resolve(u,t,a,h,beta): #resolve equacao nao-linear por método de Newton
	ctrl=2
	cont=0
	while (ctrl>0):
		cont=cont+1
		residuo=u-beta*h*(u**3-u+t)-a
		derivada=1-beta*h*(3*u**2-1)
		u1=u-residuo/derivada
		if np.abs(u1-u) <= 1e-10*(1+np.abs(u1)):
			ctrl=ctrl-1
		else:
			ctrl=2
		u=u1
#	print cont
	return u


def adams_moulton_2(h,Tmax,u1):
	itmax=int(round(Tmax/h))
	u=np.empty((itmax+1,1))
	u[0]=u1

	for i in range(itmax):
		t=i*h
		fn=f(t,u[i])
		a=u[i] + h*fn/2
		u[i+1]=resolve(u[i]+ h*fn,t+h,a,h,1/2)
	return u



u0=0
h=1e-2
Tmax=1
itmax=int(round(Tmax/h))

u=adams_moulton_2(h,Tmax,u0)
print(u[itmax])

\end{verbatim}

\fi
\end{exeresol}


\subsection*{Exercícios}

\begin{exer}
Encontre o método de Adams-Moulton para $s=0$.
\end{exer}
\begin{resp}
 \begin{equation}
  y^{(n+1)}=y^{(n)}+hf\left(t^{(n+1)},y^{(n+1)}\right)
 \end{equation}
Este esquema é equivalente ao método de Euler Implícito.
 \end{resp}



\begin{exer}
Encontre o método de Adams-Moulton para $s=1$.
\end{exer}
\begin{resp}
 \begin{equation}
  y^{(n+1)}=y^{(n)}+\frac{h}{2}\left[f\left(t^{(n+1)},u(t^{(n+1)})\right)+f\left(t^{(n)},u(t^{(n)})\right)\right]
 \end{equation}
 Este esquema é equivalente ao método trapezoidal.
 \end{resp}


 \begin{exer} Repita o Problema~\ref{exeresol:exersol_moulton_nl} usando Adams-Moulton com 3 passos e inicilizando com Runge-Kutta quarta ordem clássico.
  \end{exer}
\begin{resp}
 $0,37517345$ e  $0,37512543$.
\end{resp}


%
%
% \section{Método BDF}
% Um método de ordem $s$ com $s$ estágios é chamado de \emph{método BDF-Backward Differentiation Formula} se $\sigma (w)=b_sw^s$, onde $b_s \in \mathbb{R}$, ou seja,
% \begin{eqnarray}\label{BDF}
%   a_s u_{n+s}+ ...+ a_1 u_{n+1} + a_0u_{n} &=  h b_sf_{n+s}
% \end{eqnarray}
%
% \begin{ex}
% Mostre que o método BDF com $s=3$ é
% \begin{eqnarray}
%   u_{n+3} -\frac{18}{11} u_{n+2}+\frac{9}{11}u_{n+1}-\frac{2}{11}u^{(n)} &= \frac{6}{11}h f_{n+3}
% \end{eqnarray}
% \end{ex}
%
% \subsection*{Exercícios}
%
% \begin{exer}
% Mostre que o método BDF com $s=1$ é o método de Euler implícito.
% \end{exer}
%
% \begin{exer}
% Mostre que o método BDF com $s=2$ é
% \begin{eqnarray}
%   u_{n+2} -\frac{4}{3} u_{n+1} + \frac{1}{3}u^{(n)} &= \frac{2}{3}h f_{n+2}
% \end{eqnarray}
% \end{exer}
%
%
% \section{O método de Taylor}
% Uma maneira simples de aumentar a ordem dos métodos de Euler anteriormente descritos consiste em truncar a série de Taylor de $u(t+h)$:
% \begin{eqnarray}
%  u(t+h)=u(t) +h u'(t)+ \frac{h^2}{2!}u''(t)+\frac{h^3}{3!}u'''(t)+\ldots
% \end{eqnarray}
%
% Utilizando dois termos temos o método de Euler. Utilizando os três primeiros termos da série e substituindo $u'(t)=f(t,u)$ e $u''(t)=\frac{\partial f}{\partial t}(t,x)$ temos o \emph{método de Taylor de ordem $2$}
% \begin{eqnarray}
%    u_{n+1}=u^{(n)} +h f(t^{(n)},u^{(n)})+ \frac{h^2}{2!} \frac{\partial f}{\partial t}(t^{(n)},u^{(n)})
% \end{eqnarray}
%
%
% De forma análoga, o método de Taylor de ordem $3$ é
% \begin{eqnarray}
%    u_{n+1}=u^{(n)} +h f(t^{(n)},u^{(n)})+ \frac{h^2}{2!}\frac{\partial f}{\partial t}(t^{(n)},u^{(n)})+\frac{h^3}{3!}\frac{\partial^2 f}{\partial t^2}(t^{(n)},u^{(n)})
% \end{eqnarray}
%
% \subsection*{Exercícios resolvidos}
%
% \emconstrucao
%
% \subsection*{Exercícios}
%
% \emconstrucao
%
% \section{Estabilidade dos métodos de Taylor}
% \begin{ex}
% Prove que para um método de Taylor de ordem $p$ para a EDO \eqref{EDO4.7} temos
% \begin{eqnarray}
%   p(z)= 1 + z+ \frac{z^2}{2!} +\frac{z^3}{3!}+\ldots +\frac{z^p}{p!}
% \end{eqnarray}
% onde  $u^{(n)} = (p(z))^nu_0$ e a região de estabilidade é dada por
% \begin{eqnarray}
%  \mathcal D_{T} = \{z \in  \field{C}:  \left|p(z)\right|<1\}
% \end{eqnarray}
%
% Trace as regiões de estabilidade para o método de Taylor para $p=1,\ldots ,6$ no mesmo gráfico.
% \end{ex}
%
% % \begin{figure}
% % \begin{center}
% %   \includegraphics[width=8cm]{RegiaoTaylor.eps}\\
% %   \caption{Região de estabilidade paras os métodos de Taylor de ordem $1,\ldots ,4$ (interior as curvas). A curva mais interna é para $p=1$}\label{RegiaoTaylor}
% % \end{center}
% % \end{figure}
%
%
% \begin{ex}
% Aproxime a solução do problema de valor inicial
% \begin{eqnarray}
%    \frac{du}{dt} &=\sin{t}\\
%             u(0) &= 1
% \end{eqnarray}
% para  $t\in [0,10]$.
%
% \begin{enumerate}
% \item [a.] Trace a solução para $h=0,16$, $0,08$, $0,04$, $0,02$ e $0,01$ para o método de Taylor de ordem $1$, $2$ e $3$. (Trace todos de ordem $1$ no mesmo gráfico, ordem $2$ em outro gráfico e ordem $3$ outro gráfico separado.)
%
% \item [b.] Utilizando a solução exata, trace um gráfico do erro em escala logar\'itmica.
% Comente os resultados (novamente, em cada gráfico separado para cada método repita os valores acima)
%
% \item [c.] Fixe agora o valor $h=0,02$ e trace no mesmo gráfico uma curva para cada método.
%
% \item [d.] Trace em um gráfico o erro em $t=10$ para cada um dos métodos (uma curva para cada ordem) a medida que $h$ diminui. (Use escala \verb#loglog#)
% \end{enumerate}
% \end{ex}

\section{Método de Adams-Moulton para sistemas lineares}
Esquemas implícitos como o de Adams-Moulton apresentam a dificuldade adicional de necessitar do valor de $f(t^{(n+1)},u^{(n+1)})$ para calcular o valor de $u^{(n+1)}$. Pelo menos para sistemas lineares, o método pode ser explicitado. Seja o seguinte problema de valor inicial linear:
\begin{eqnarray}
 u'(t)&=& Au(t) + g(t),\\
 u(t^{(1)}) &=& a.
\end{eqnarray}
Onde $u(t)$ é um vetor de n entradas e $A$ é uma matriz $n\times n$.

Considere agora o esquema de Adams-Moulton dado na Equação~\eqref{eq:mam} com $f(t,u)=Au+g(t)$:
\begin{eqnarray}
  u^{(n+s)}  &=& u^{(n+s-1)}  + h\sum_{j=0}^{s}\beta_j \left[Au^{(n+j)} + g(t^{(n+j)})\right]
\end{eqnarray}
o que pode ser escrito como:
\begin{eqnarray}
  \left(I_d-h\beta_s A\right)u^{(n+s)}  &=& u^{(n+s-1)}  + h\sum_{j=0}^{s-1}\beta_j \left[Au^{(n+j)}+ g(t^{(n+j)})\right]\nonumber\\
  &+& h\beta_s g(t^{(n+s)}) \label{eq:adams:lineares}
\end{eqnarray}
onde $I_d$ é matriz identidade $n\times n$. O sistema linear envolvido em \eqref{eq:adams:lineares} pode ser resolvido sempre que $I_d-h\beta_s A$ for inversível, o que sempre acontece quando $h$ é suficientemente pequeno.

\subsection*{Exercícios resolvidos}

\construirExeresol

\subsection*{Exercícios}

\construirExer

\section{Estratégia preditor-corretor}\index{preditor-corretor}
Esquemas implícitos como o de Adams-Moulton (Seção~\ref{sec:Adams_Moulton}) e o de Runge-Kutta (Seção~\ref{sec:sec_IRK}), embora úteis para resolver problemas rígidos (ver Seção~\ref{sec:stiff}), apresentam a dificuldade de necessitar do valor de $f(t^{(n+1)},u^{(n+1)})$ para calcular o valor de $u^{(n+1)}$, exigindo a solução de uma equação algébrica a cada passo. Uma forma de evitar a solução completa da equação implícita a cada passo consiste em estimar o valor futuro por um método explícito e usar essa estimativa na etapa de correção. Essa abordagem é chamada de \emph{estratégia preditor-corretor}.

Os métodos do tipo preditor-corretor empregam um esquema explícito para \emph{predizer} o valor de $u^{(n+1)}$ e, depois, um método implícito para recalcular, isto é, \emph{corrigir} $u^{(n+1)}$.

\begin{ex}
Considere o método de Euler implícito aplicado ao problema
\begin{eqnarray}
u'(t)&=&f(t,u(t)),\\
u(t^{(1)})&=&a.
\end{eqnarray}
A recorrência implícita é
\begin{eqnarray}
u^{(n+1)}&=&u^{(n)}+hf(t^{(n+1)},u^{(n+1)}).
\end{eqnarray}
Podemos usar o método de Euler explícito como preditor:
\begin{eqnarray}
\tilde{u}^{(n+1)}&=&u^{(n)}+hf(t^{(n)},u^{(n)}).
\end{eqnarray}
Substituindo o valor predito na avaliação implícita, obtemos uma única correção:
\begin{eqnarray}
u^{(n+1)}&=&u^{(n)}+hf(t^{(n+1)},\tilde{u}^{(n+1)}).
\end{eqnarray}
Esse esquema é um método preditor-corretor explícito obtido a partir de Euler implícito. Ele não coincide com o método de Euler melhorado da Seção~\ref{sec:sec_euler_mod}, pois este último emprega a média das inclinações no início e no fim do passo.
\end{ex}


\begin{ex} Considere o método de trapezoidal (ver \ref{sec:trapezoidal}) aplicado para resolver o problema de valor inicial
\begin{eqnarray}
  u'(t)  &=& f(t,u(t)) \\
  u(t^{(1)}) &=& a
\end{eqnarray}
cujo processo iterativo é dado por
\begin{eqnarray}
u^{(n+1)}=u^{(n)} + \frac{h}{2} \left[f(t^{(n)},u^{(n)})+f(t^{(n+1)},u^{(n+1)})\right].
\end{eqnarray}
Agora aplicamos o método de Euler (ver \ref{sec:euler}) para predizer $u^{(n+1)}$:
\begin{eqnarray}
\tilde{u}^{(n+1)}=u^{(n)} + h f(t^{(n)},u^{(n)}).
\end{eqnarray}
E agora, retornamos ao método trapezoidal para obter:
\begin{eqnarray}
\tilde{u}^{(n+1)}&=&u^{(n)} + h f(t^{(n)},u^{(n)}),\\
u^{(n+1)}&=&u^{(n)} + \frac{h}{2} \left[f(t^{(n)},u^{(n)})+f(t^{(n+1)},\tilde{u}^{(n+1)})\right].
\end{eqnarray}
\end{ex}

\begin{ex} Considere o método de Adams-Moulton de segunda ordem (ver \ref{sec:Adams_Moulton}) aplicado para resolver  o problema de valor inicial
\begin{eqnarray}
  u'(t)  &=& f(t,u(t)) \\
  u(t^{(1)}) &=& a
\end{eqnarray}
cujo processo iterativo é dado por
\begin{eqnarray}
u^{(n+1)}=u^{(n)} + \frac{h}{2} \left[f(t^{(n)},u^{(n)})+f(t^{(n+1)},u^{(n+1)})\right].
\end{eqnarray}
Agora aplicamos o método de Adams-Bashforth de segunda ordem (ver \ref{sec:Adams_Bashforth}) para predizer $u^{(n+1)}$:
\begin{eqnarray}
\tilde{u}^{(n+1)}=u^{(n)} + \frac{h}{2} \left[-f(t^{(n-1)},u^{(n-1)})+3f(t^{(n)},u^{(n)})\right].
\end{eqnarray}
Assim, obtemos o seguinte método:
\begin{eqnarray}
\tilde{u}^{(n+1)}&=&u^{(n)} + \frac{h}{2} \left[-f(t^{(n-1)},u^{(n-1)})+3f(t^{(n)},u^{(n)})\right],\\
u^{(n+1)}&=&u^{(n)} + \frac{h}{2} \left[f(t^{(n)},u^{(n)})+f(t^{(n+1)},\tilde{u}^{(n+1)})\right].
\end{eqnarray}
\end{ex}

\subsection{Exercícios resolvidos}

\construirExeresol

\subsection*{Exercícios}

\begin{exer} Construa o esquema preditor corretor combinando Adams-Moulton de quarta ordem e Adams-Bashforth de quarta ordem.
\end{exer}
\begin{resp}
\begin{eqnarray}
\tilde{u}^{(n+1)}&=&u^{(n)} + \frac{h}{24} \left[-9f(t^{(n-3)},u^{(n-3)}) + 37f(t^{(n-2)},u^{(n-2)}) - 59 f(t^{(n-1)},u^{(n-1)})  + 55f(t^{(n)},u^{(n)})\right],\\
u^{(n+1)}&=&u^{(n)} + \frac{h}{24} \left[f(t^{(n-2)},u^{(n-2)})-5f(t^{(n-1)},u^{(n-1)})+19f(t^{(n)},u^{(n)})+9 f(t^{(n+1)},\tilde{u}^{(n+1)})\right].
\end{eqnarray}
 \end{resp}

 \begin{exer} Seja o problema de valor inicial dado por:
\begin{eqnarray}
  u'(t)  &=& \sqrt{u(t)+1} \\
  u(0) &=& 0
\end{eqnarray} Resolva numericamente esse problema pelo método de Adams-Bashforth de segunda ordem e pelo método preditor corretor combinando Adams-Bashforth de segunda ordem com Adams-Moulton de segunda ordem. Compare a solução obtida para $t=10$ com a solução exata dada por:
\begin{equation} u(t)=\frac{t^2}{4}+t. \end{equation}
 Inicialize os métodos empregando Runge-Kutta de segunda ordem.
\end{exer}
\begin{resp}
 Adams-Bashforth: 34,99965176,~~ Preditor-corretor: 34,99965949,~~  Exato: 35
\end{resp}

\section{Problemas rígidos}\index{problema!rígido}\label{sec:stiff}

Em alguns problemas de valor inicial, a escolha do passo é limitada mais pela estabilidade do método numérico do que pela precisão desejada. Esse comportamento ocorre com frequência quando a solução contém componentes que evoluem em escalas de tempo muito diferentes. Tais problemas são chamados de \emph{problemas rígidos}.

Não existe uma única definição de rigidez adequada a todas as situações. Uma caracterização operacional importante é a seguinte: um problema é rígido quando um método explícito estável exige passos muito menores do que aqueles que seriam necessários apenas para representar a variação da solução de interesse.

A equação teste
\begin{eqnarray}\label{eq:teste_estabilidade}
u'(t)&=&\lambda u(t)
\end{eqnarray}
é útil para estudar esse fenômeno. Quando $\mathop{\rm Re}(\lambda)<0$, sua solução
\begin{equation}
u(t)=u(0)e^{\lambda t}
\end{equation}
decai para zero.

Aplicando o método de Euler explícito, obtemos
\begin{eqnarray}
u^{(n+1)}&=&(1+h\lambda)u^{(n)}.
\end{eqnarray}
Definindo $z=h\lambda$, o fator de amplificação é
\begin{equation}
R(z)=1+z.
\end{equation}
Para que a solução numérica da equação teste também decaia, é necessário que
\begin{equation}
|1+z|<1.
\end{equation}
No caso $\lambda<0$ real, essa condição equivale a
\begin{equation}
0<h<\frac{2}{|\lambda|}.
\end{equation}
Assim, um modo com $|\lambda|$ muito grande força o uso de um passo muito pequeno, mesmo quando esse modo já se tornou praticamente desprezível na solução.

\begin{ex}
Considere o sistema
\begin{eqnarray}
u_1'(t)&=&-u_1(t),\\
u_2'(t)&=&-1000u_2(t),
\end{eqnarray}
com $u_1(0)=u_2(0)=1$. A solução exata é
\begin{eqnarray}
u_1(t)&=&e^{-t},\\
u_2(t)&=&e^{-1000t}.
\end{eqnarray}
A segunda componente desaparece em uma escala de tempo da ordem de $10^{-3}$, enquanto a primeira varia em uma escala da ordem de $1$. Entretanto, o método de Euler explícito exige
\begin{equation}
h<2\times 10^{-3}
\end{equation}
para tratar de forma estável a segunda componente. Portanto, mesmo depois que $u_2(t)$ se tornou muito pequena, o passo continua restrito pela escala rápida.
\end{ex}

Métodos implícitos podem ser mais adequados nesse contexto. Para Euler implícito aplicado à Equação~\eqref{eq:teste_estabilidade},
\begin{eqnarray}
u^{(n+1)}&=&\frac{1}{1-z}u^{(n)},
\end{eqnarray}
de modo que
\begin{equation}
R(z)=\frac{1}{1-z}.
\end{equation}
Se $\mathop{\rm Re}(z)<0$, então $|R(z)|<1$ para qualquer valor de $h$. Dizemos, nesse caso, que o método é \emph{$A$-estável}. Além disso,
\begin{equation}
\lim_{z\to-\infty}R(z)=0,
\end{equation}
propriedade chamada de \emph{$L$-estabilidade}. Ela faz com que componentes muito rápidas sejam fortemente amortecidas numericamente.

O método trapezoidal também é $A$-estável, pois
\begin{equation}
R(z)=\frac{1+z/2}{1-z/2}.
\end{equation}
Entretanto,
\begin{equation}
\lim_{z\to-\infty}R(z)=-1.
\end{equation}
Assim, modos muito rígidos não são amortecidos pelo método trapezoidal; eles podem aparecer na solução numérica como oscilações alternadas de pequena ou moderada amplitude. Esse exemplo mostra que estabilidade absoluta e amortecimento de modos rígidos são propriedades relacionadas, mas não idênticas.

Na prática, a rigidez também afeta a escolha de algoritmos adaptativos. Métodos explícitos podem reduzir o passo repetidamente por razões de estabilidade, enquanto métodos implícitos toleram passos maiores ao custo de resolver sistemas algébricos em cada passo. A escolha entre esses métodos depende do custo de cada passo, da precisão desejada e da estrutura do problema.

\section{Verificação, validação e hierarquia de ``benchmarks''}\index{verificação}\index{validação}\index{Benchmark}

Ao aplicar uma metodologia numérica, é necessário avaliar se o resultado computado é confiável. Convém distinguir dois conceitos. A terminologia adotada aqui segue a literatura de verificação e validação em computação científica \cite{AIAA1998VV,ASMEVV10,AIAA2024Code,OberkampfRoy2025}.

\begin{itemize}
\item \emph{Verificação} pergunta se as equações matemáticas foram resolvidas corretamente. Inclui a detecção de erros de implementação, a análise do erro de discretização, o controle de erros algébricos e a investigação dos efeitos de arredondamento.
\item \emph{Validação} pergunta se o modelo matemático representa adequadamente o fenômeno que se deseja descrever. Em problemas provenientes de aplicações, essa etapa normalmente envolve comparação com dados experimentais ou observacionais.
\end{itemize}

Neste capítulo, como as equações diferenciais são tomadas como dadas, a maior parte dos testes é de \emph{verificação numérica}, e não de validação do modelo. Também é útil distinguir \emph{verificação do código}, que procura estabelecer se a implementação corresponde ao método matemático pretendido, e \emph{verificação da solução}, que procura quantificar os erros presentes em uma solução calculada para um problema particular \cite{Roy2005Verification,OberkampfRoy2025}.

Uma solução de referência de alta confiança usada para comparação será chamada de {\it benchmark}. A literatura de verificação não estabelece uma ordem matemática absoluta entre todos os possíveis {\it benchmarks}, mas a prática consolidou uma hierarquia operacional de problemas de teste. Os guias da AIAA e da ASME e trabalhos posteriores organizam os problemas, em ordem aproximadamente decrescente de confiança na solução de referência, como: soluções analíticas exatas, incluindo soluções manufaturadas; soluções semi-analíticas; e soluções numéricas de alta precisão \cite{AIAA1998VV,ASMEVV10,Anderson2007VV,OberkampfTrucano2008,AIAA2024Code}. Essa hierarquia deve ser entendida como um princípio prático. Uma fórmula exata difícil de avaliar numericamente pode ser menos útil que um {\it benchmark} numérico acompanhado por um controle de erro muito rigoroso. Além disso, um bom problema de verificação deve exercitar suficientemente os termos e mecanismos presentes no método ou no código que se deseja testar \cite{Roache2002MMS,KnuppSalari2003}.

Com essa ressalva, podemos organizar as formas de verificação da seguinte maneira.

\begin{enumerate}
\item \textbf{Soluções analíticas exatas e soluções manufaturadas.} Constituem, em geral, a referência mais forte para verificação de código \cite{Roache2002MMS,OberkampfRoy2025}. Quando uma solução exata está disponível, o erro pode ser calculado diretamente e a ordem observada do método pode ser comparada com a ordem teórica. Quando não existe uma solução exata suficientemente rica para o problema original, pode-se escolher uma função suave e construir os termos fontes e as condições de contorno de modo que essa função seja solução exata do problema modificado. Essa técnica é conhecida como \emph{método das soluções manufaturadas}.

Dentro da classe das soluções exatas existe ainda uma gradação prática associada à dificuldade de avaliar a solução de referência. Essa gradação não constitui uma nova hierarquia matemática, mas é relevante para o cálculo do {\it benchmark}:

\subitem Expressões aritméticas ou algébricas, que podem ser avaliadas usando um número finito de operações elementares, como $u(t)=\frac{t^2+1}{3t-4}$ ou $u(t)=t^{3/2}+\sqrt{2}$.

\subitem Expressões em forma fechada envolvendo funções elementares, como exponenciais, logaritmos e funções trigonométricas. Por exemplo, $u(t)=e^{-t}\sin(t)$.

\subitem Expressões envolvendo funções especiais. Essas expressões continuam representando soluções matematicamente exatas, mas sua avaliação depende de algoritmos numéricos para as funções especiais. É necessário, portanto, que o erro introduzido nessa avaliação seja muito menor que o erro que se deseja medir no método testado.

\subitem Representações exatas por séries ou outras expansões infinitas também podem ser usadas, desde que o erro de truncamento da representação seja controlado. Em certos casos, uma representação formalmente exata pode ser inconveniente como {\it benchmark} se sua avaliação numérica for tão difícil quanto a solução do problema original \cite{Roache2002MMS}.

\item \textbf{Soluções semi-analíticas e reformulações analíticas.} O problema pode ser reduzido, por análise matemática, a uma quadratura, a uma equação implícita, a um sistema de equações ordinárias mais simples, a uma série convergente ou a uma expansão assintótica. A etapa final ainda exige cálculo numérico, mas esse cálculo é realizado por uma rota diferente daquela empregada pelo método que se deseja testar. Na terminologia usada em guias de verificação, essa classe aparece frequentemente como \emph{solução semi-analítica} \cite{AIAA1998VV,ASMEVV10}. Para que a comparação seja útil, o erro introduzido na avaliação semi-analítica deve ser controlado e suficientemente menor que o erro da solução em teste.

\item \textbf{{\it Benchmarks} numéricos de alta precisão.} Quando não existe uma solução analítica ou semi-analítica adequada, uma solução numérica pode servir de referência. Entretanto, não basta que ela tenha sido calculada com uma malha aparentemente muito fina. Um {\it benchmark} numérico forte deve ser produzido por um código previamente verificado e bem documentado, acompanhado por estimativa do erro numérico e, sempre que possível, confirmado de forma independente por outro método ou outro código \cite{OberkampfTrucano2008}. O uso de precisão estendida e de métodos de natureza diferente aumenta a independência da referência.

\item \textbf{Comparação entre métodos e implementações independentes.} Quando não se dispõe de um {\it benchmark} propriamente dito, a concordância entre métodos distintos constitui evidência útil. Essa evidência é mais forte quando os métodos não compartilham a mesma formulação, o mesmo código, a mesma discretização ou as mesmas hipóteses de implementação. A concordância, por si só, não transforma nenhum dos resultados em solução de referência.

\item \textbf{Convergência numérica por refinamento.} Calculam-se soluções com passos sucessivamente menores e verifica-se se as diferenças diminuem com a ordem prevista. Esse teste é fundamental para a verificação, em especial quando associado ao teste de ordem de precisão \cite{Roy2005Verification,KnuppSalari2003}, mas, isoladamente, não fornece uma referência externa. Uma sequência pode convergir para uma solução incorreta devido a um erro sistemático na formulação ou na implementação.

\item \textbf{Resíduos, invariantes e propriedades qualitativas.} Conservação de massa ou energia, positividade, monotonicidade, simetrias, limites assintóticos e resíduos da equação fornecem testes complementares. Essas propriedades não ocupam necessariamente um nível específico da hierarquia acima: elas constituem informações independentes que podem fortalecer qualquer um dos testes anteriores.
\end{enumerate}

A hierarquia acima pode ser resumida esquematicamente por
\begin{equation}
\begin{array}{c}
\hbox{solução analítica exata ou manufaturada}\\
\Downarrow\\
\hbox{solução semi-analítica ou reformulação independente}\\
\Downarrow\\
\hbox{{\it benchmark} numérico de alta precisão}\\
\Downarrow\\
\hbox{comparação entre métodos independentes}\\
\Downarrow\\
\hbox{autoconvergência por refinamento}.
\end{array}
\end{equation}
As setas indicam uma redução típica da informação independente disponível sobre a solução de referência e não uma relação matemática de ordem. Na prática, devem-se combinar tantos testes quanto forem razoavelmente possíveis.

\begin{ex}[Expressão analítica]
O problema
\begin{eqnarray}
u'(t)&=&-0,5u(t)+2+t,\\
u(0)&=&8
\end{eqnarray}
possui solução exata
\begin{equation}
u(t)=2t+8e^{-t/2}.
\end{equation}
Esse problema permite calcular diretamente o erro da aproximação e verificar a ordem observada do método.
\end{ex}

\begin{ex}[Solução manufaturada]
Considere a função escolhida a priori
\begin{equation}
u(t)=e^{-t}\sin(t).
\end{equation}
Podemos construir o problema
\begin{eqnarray}
u'(t)&=&u(t)^2+g(t),\\
u(0)&=&0,
\end{eqnarray}
onde
\begin{equation}
g(t)=e^{-t}(\cos(t)-\sin(t))-e^{-2t}\sin^2(t).
\end{equation}
Por construção, $u(t)=e^{-t}\sin(t)$ é solução exata. O objetivo desse problema não é representar um modelo físico particular, mas testar se a implementação reproduz a ordem e o comportamento esperados do método numérico. Essa é a ideia central do método das soluções manufaturadas \cite{Roache2002MMS}.
\end{ex}

\begin{ex}[Expressão envolvendo função especial]
A solução do problema
\begin{eqnarray}
u'(t)&=&-u^3(t)+u^2(t),\\
u(0)&=&\frac{1}{2}
\end{eqnarray}
pode ser escrita como
\begin{equation}
u(t)=\frac{1}{1+W(e^{1-t})},
\end{equation}
onde $W$ é a função de Lambert, definida por
\begin{equation}
W(x)e^{W(x)}=x.
\end{equation}
Embora $W$ não seja uma função elementar, bibliotecas de funções especiais permitem calcular essa solução com alta precisão. Para usá-la como {\it benchmark}, a precisão da avaliação de $W$ deve ser superior à precisão que se deseja testar na solução numérica.
\end{ex}

\begin{ex}[Reformulação analítica implícita]
Considere
\begin{eqnarray}
u'(t)&=&u^3(t)+u^2(t)+u(t)+1,\\
u(0)&=&0.
\end{eqnarray}
Por separação de variáveis e integração, obtém-se
\begin{equation}
\ln\left(\frac{(u(t)+1)^2}{u(t)^2+1}\right)+2\arctan(u(t))=4t.
\end{equation}
Essa relação não fornece $u$ explicitamente em função de $t$, mas fornece $t$ explicitamente em função de $u$. Por exemplo,
\begin{equation}
u=1
\end{equation}
corresponde a
\begin{equation}
t=\frac{\ln(2)}{4}+\frac{\pi}{8}\approx0,5659858768387104.
\end{equation}
Além disso, tomando $u\to+\infty$, obtemos
\begin{equation}
t\to\frac{\pi}{4}-.
\end{equation}
Assim, o mesmo problema fornece referências exatas para o valor da solução em um ponto e para o tempo de explosão. Essas informações são obtidas por uma rota independente do integrador de EDO que se deseja verificar.
\end{ex}

\begin{ex}[Reformulação por quadratura]
Considere
\begin{eqnarray}
u'(t)&=&[\cos(u(t))+u(t)](1+\cos(t)),\\
u(0)&=&0.
\end{eqnarray}
A separação de variáveis fornece
\begin{equation}
\int_0^{u(t)}\frac{d\tau}{\cos(\tau)+\tau}=t+\sin(t).
\end{equation}
Podemos escolher um valor de $u$ e calcular o tempo correspondente por um procedimento independente. Para $u=100$,
\begin{equation}
\int_0^{100}\frac{d\tau}{\cos(\tau)+\tau}
\approx5,574304717298400.
\end{equation}
Resolvendo
\begin{equation}
t+\sin(t)=5,574304717298400,
\end{equation}
obtemos
\begin{equation}
t_*\approx5,924938036503083.
\end{equation}
Portanto,
\begin{equation}
u(t_*)=100
\end{equation}
fornece um {\it benchmark} semi-analítico. Observe que a quadratura e a solução da equação escalar são problemas numericamente diferentes da integração direta da EDO.

A tabela a seguir ilustra a comparação desse {\it benchmark} com quatro métodos de quarta ordem. Para evitar introduzir um erro adicional pela escolha do tempo final, tomamos $h=t_*/N$, de forma que todos os métodos sejam avaliados exatamente em $t=t_*$. Os métodos de Adams são inicializados com Runge-Kutta clássico.

\begin{center}
\begin{tabular}{|c|c|c|c|}
\hline
&$N=60$&$N=600$&$N=6000$\\
\hline
Runge-Kutta 4&100,5019292&99,99992527&99,99999999\\
\hline
Adams-Bashforth 4&99,74962359&99,98728377&100,00000458\\
\hline
Preditor-corretor 4&100,01697360&99,99787111&99,99999971\\
\hline
Adams-Moulton 4&99,94262050&99,99736562&99,99999969\\
\hline
{\it Benchmark}&100&100&100\\
\hline
\end{tabular}
\end{center}

A concordância progressiva dos quatro métodos com o valor de referência reforça a evidência fornecida pela reformulação analítica. A tabela também ilustra por que uma simples comparação entre métodos é mais fraca que a comparação com uma referência construída independentemente: os quatro métodos poderiam concordar entre si sem que isso, isoladamente, determinasse o valor correto.
\end{ex}


\section{Convergência, consistência e estabilidade}

A análise de erro de um método numérico envolve três conceitos relacionados: consistência, estabilidade e convergência.

\begin{defn}
Um método de passo simples escrito na forma
\begin{equation}
u^{(n+1)}=u^{(n)}+h\Phi(t^{(n)},u^{(n)},h)
\end{equation}
é \emph{consistente} com a equação
\begin{equation}
u'=f(t,u)
\end{equation}
se
\begin{equation}
\lim_{h\to0}\Phi(t,u,h)=f(t,u).
\end{equation}
\end{defn}

A consistência garante que um único passo reproduz a equação diferencial no limite $h\to0$. Ela não é, sozinha, suficiente para garantir que os erros introduzidos ao longo de muitos passos permaneçam controlados.

\begin{defn}
Um método é \emph{convergente} em $[t^{(1)},T]$ se, para toda solução suficientemente regular do problema de valor inicial,
\begin{equation}
\max_{t^{(n)}\leq T}|u(t^{(n)})-u^{(n)}|\longrightarrow0
\end{equation}
quando $h\to0$.
\end{defn}

Para métodos de passo simples cuja função de incremento é Lipschitz contínua na variável dependente, a consistência e uma estimativa adequada de estabilidade permitem provar a convergência. Em termos informais, a consistência controla o erro produzido em cada passo e a estabilidade controla sua propagação.

Uma forma diferente de estabilidade, importante na prática, é a \emph{estabilidade absoluta}. Aplicamos o método à equação teste
\begin{equation}
u'=\lambda u
\end{equation}
e escrevemos o resultado na forma
\begin{equation}
u^{(n+1)}=R(z)u^{(n)},~~~z=h\lambda.
\end{equation}
A função $R(z)$ é chamada de \emph{função de estabilidade}. A região de estabilidade absoluta é
\begin{equation}
\mathcal{S}=\{z\in\mathbb{C}:|R(z)|\leq1\}.
\end{equation}

Para alguns métodos estudados neste capítulo, temos:
\begin{center}
\begin{tabular}{|c|c|}
\hline
Método & $R(z)$\\
\hline
Euler explícito & $1+z$\\
\hline
Euler melhorado & $1+z+\dfrac{z^2}{2}$\\
\hline
Euler implícito & $\dfrac{1}{1-z}$\\
\hline
Trapezoidal & $\dfrac{1+z/2}{1-z/2}$\\
\hline
Runge-Kutta clássico de ordem 4 &
$1+z+\dfrac{z^2}{2}+\dfrac{z^3}{6}+\dfrac{z^4}{24}$\\
\hline
\end{tabular}
\end{center}

A estabilidade absoluta não deve ser confundida com a estabilidade de um método de passo múltiplo em relação a perturbações nos valores iniciais. Para um método linear de passo múltiplo,
\begin{equation}
\sum_{j=0}^s\alpha_j u^{(n+j)}
=
h\sum_{j=0}^s\beta_j f^{(n+j)},
\end{equation}
define-se o polinômio
\begin{equation}
\rho(\xi)=\sum_{j=0}^s\alpha_j\xi^j.
\end{equation}
O método é \emph{zero-estável} quando todas as raízes de $\rho$ satisfazem
\begin{equation}
|\xi|\leq1
\end{equation}
e toda raiz com $|\xi|=1$ é simples. Para métodos lineares de passo múltiplo, sob as hipóteses usuais de regularidade, vale o princípio fundamental:
\begin{equation}
\hbox{consistência + zero-estabilidade}
\quad\Longleftrightarrow\quad
\hbox{convergência}.
\end{equation}

\subsection*{Ordem observada}

Quando não conhecemos a solução exata, ainda podemos estimar numericamente a ordem do método. Suponha que, em um ponto fixo, o erro dominante tenha a forma
\begin{equation}
u_h-u=C h^p+O(h^{p+1}).
\end{equation}
Usando passos $h$, $h/2$ e $h/4$, obtemos a estimativa
\begin{equation}
p_{\rm obs}\approx
\log_2\left(
\frac{|u_h-u_{h/2}|}{|u_{h/2}-u_{h/4}|}
\right).
\end{equation}
Para sistemas, o valor absoluto pode ser substituído por uma norma.

Essa estimativa só é confiável quando os passos já estão na região assintótica em que o termo $Ch^p$ domina o erro. Para passos muito grandes, o comportamento assintótico pode ainda não ter aparecido; para passos muito pequenos, erros de arredondamento ou tolerâncias internas podem dominar.

\begin{obs}
A convergência observada sob refinamento é uma evidência importante, mas deve ser interpretada junto com outros testes. Um erro de implementação que permaneça praticamente inalterado com $h$, por exemplo, pode fazer uma sequência de resultados parecer convergente para um valor incorreto.
\end{obs}

\section{Exercícios finais}


\begin{exer} Considere o problema de valor inicial dado por
\begin{eqnarray}
\frac{d u(t)}{dt} &=& -u(t) + e^{-t} \\
u(0)&=&0
\end{eqnarray}
Resolva analiticamente este problema usando as técnicas elementares de equações diferenciais ordinárias. A seguir encontre aproximações numéricas usando os métodos de Euler, Euler modificado, Runge-Kutta clássico e Adams-Bashforth de ordem 4 conforme pedido nos itens.
\begin{itemize}
\item[a)] Use passo $h=0,1$. Construa uma tabela apresentando valores com 7 algarismos significativos para comparar a solução analítica com as aproximações numéricas produzidas pelos métodos sugeridos. Construa também uma tabela para o erro absoluto obtido por cada método numérico em relação à solução analítica. Nesta última tabela, expresse o erro com 2 algarismos significativos em formato científico. Dica: $format('e',8)$ para a segunda tabela.
\begin{center}
\begin{tabular}{|c|c|c|c|c|c|}
\hline
&0,5&1,0&1,5&2,0&2,5\\
\hline
Analítico&&&&&\\
\hline
Euler&&&&&\\
\hline
Euler modificado&&&&&\\
\hline
Runge-Kutta clássico&&&&&\\
\hline
Adams-Bashforth ordem 4&&&&&\\
\hline
\end{tabular}
\end{center}

\begin{center}
\begin{tabular}{|c|c|c|c|c|c|}
\hline
&0,5&1,0&1,5&2,0&2,5\\
\hline
Euler&&&&&\\
\hline
Euler modificado&&&&&\\
\hline
Runge-Kutta clássico&&&&&\\
\hline
Adams-Bashforth ordem 4&&&&&\\
\hline
\end{tabular}
\end{center}

\item[b)] Calcule o valor produzido por cada um desses método para $u(1)$ com passo $h=0,1$, $h=0,05$, $h=0,01$, $h=0,005$ e $h=0,001$. Complete a tabela com os valores para o erro absoluto encontrado.
\begin{center}
\begin{tabular}{|c|c|c|c|c|c|}
\hline
&0,1&0,05&0,01&0,005&0,001\\
\hline
Euler&&&&&\\
\hline
Euler modificado&&&&&   \\
\hline
Runge-Kutta clássico&&&&&\\
\hline
Adams-Bashforth ordem 4&&&&&\\
\hline
\end{tabular}
\end{center}

\end{itemize}

\end{exer}


\begin{resp}

\begin{center}
\begin{tabular}{|c|c|c|c|c|c|}
\hline
&0,5&1,0&1,5&2,0&2,5\\
\hline
Analítico&  0,3032653 &   0,3678794  &  0,3346952  &  0,2706706 &   0,2052125  \\
\hline
Euler& 0,3315955 &   0,3969266 &   0,3563684 &   0,2844209  &  0,2128243\\
\hline
Euler modificado &0,3025634 &   0,3671929 &   0,3342207 &   0,2704083  &  0,2051058 \\
\hline
Runge-Kutta clássico& 0,3032649  &  0,3678790  &  0,3346949  &  0,2706703  &  0,2052124\\
\hline
Adams-Bashforth ordem 4& 0,3032421  &  0,3678319 &   0,3346486  &  0,2706329  &  0,2051848  \\
\hline
\end{tabular}
\end{center}


\begin{center}
\begin{tabular}{|c|c|c|c|c|c|}
\hline
&0,5&1,0&1,5&2,0&2,5\\
\hline
Euler& 2,8e-2  &  2,9e-2  &  2,2e-2  &  1,4e-2 &   7,6e-3\\
\hline
Euler modificado& 7,0e-4  &  6,9e-4   & 4,7e-4 &   2,6e-4 &   1,1e-4\\
\hline
Runge-Kutta clássico& 4,6e-7 &   4,7e-7    &3,5e-7  &  2,2e-7 &   1,2e-7\\
\hline
Adams-Bashforth ordem 4&  2,3e-5 &   4,8e-5  &  4,7e-5  &  3,8e-5  &  2,8e-5 \\
\hline
\end{tabular}
\end{center}

\begin{center}
\begin{tabular}{|c|c|c|c|c|c|}
\hline
&0,1&0,05&0,01&0,005&0,001\\
\hline
Euler&2,9e-2&1,4e-2&2,8e-3&1,4e-3&2,8e-4\\
\hline
Euler modificado&6,9e-4&1,6e-4&6,2e-6&1,5e-6&6,1e-8\\
\hline
Runge-Kutta clássico&4,7e-7&2,8e-8&4,3e-11&2,6e-12&4,6e-15\\
\hline
Adams-Bashforth ordem 4&4,8e-5&3,3e-6&5,7e-9&3,6e-10&5,8e-13  \\
\hline
\end{tabular}
\end{center}

\end{resp}

\begin{exer} Considere o seguinte modelo para o crescimento de uma colônia de bactérias, baseado na equação logística (ver \eqref{eq:logistica})
\begin{equation} u'(t)=\alpha u(t) \left(A-u(t)\right) \end{equation}
onde $u(t)$ indica a densidade de bactérias em unidades arbitrárias na colônia e $\alpha$ e $A$ são constantes positivas.
Pergunta-se: % Solução $u(t)=\frac{Au_0}{(A-u_0)e^{-A\alpha t}+u_0}$
\begin{itemize}
\item[a)] Se $A=10$ e $\alpha=1$ e $u(0)=1$, use métodos numéricos para obter aproximação para $u(t)$ em $t=5\cdot 10^{-2}$, $t=10^{-1}$, $t=5\cdot 10^{-1}$ e $t=1$.
\item[b)] Se $A=10$ e $\alpha=1$ e $u(0)=1$, use métodos numéricos para obter tempo necessário para que a população dobre?
\item[c)] Se $A=10$ e $\alpha=1$ e $u(0)=4$, use métodos numéricos para obter tempo necessário para que a população dobre?
\end{itemize}
\end{exer}
\begin{resp}
\begin{itemize}
 \item [a)] $1,548280989603$,   $2,319693166841$,    $9,42825618574$ e  $9,995915675174$.
 \item [b)] $0,081093021622$.
 \item [c)] $0,179175946923$.
 \end{itemize}

Obs: A solução analitica do problema de valor inicial é dada por:
 \begin{equation} u(t)=\frac{Au_0}{(A-u_0)e^{-A\alpha t}+u_0} \end{equation}
 Os valores exatos para os itens b e c são:$\frac{1}{10}\ln\left(\frac{9}{4}\right)$ e $\frac{1}{10}\ln\left(6\right)$.

\end{resp}

\begin{exer} Considere o seguinte modelo para a evolução da velocidade de um objeto em queda:
\begin{equation} v'=g-\alpha v^2 \end{equation}
Sabendo que $g=9,8$ e $\alpha=10^{-2}$ e $v(0)=0$. Pede-se a velocidade ao tocar o solo e o instante quando isto acontece, dado que a altura inicial era 100.
\end{exer}
\begin{resp}
A velocidade exata ao atingir o solo é
\begin{equation}
v=\sqrt{\frac{g}{\alpha}\left(1-e^{-200\alpha}\right)}
\approx29,109644835142.
\end{equation}
O instante correspondente é
\begin{equation}
t=\frac{1}{\sqrt{g\alpha}}\cosh^{-1}\left(e^{100\alpha}\right)
=\frac{1}{\sqrt{g\alpha}}\tanh^{-1}\left(\sqrt{1-e^{-200\alpha}}\right)
\approx5,294544041566.
\end{equation}
\end{resp}



\begin{exer} Considere o seguinte modelo para o oscilador não linear de Van der Pol:
\begin{equation} u''(t) - \alpha (A-u(t)^2)u'(t) + w_0^2u(t)=0 \end{equation}
onde $A$, $\alpha$ e $w_0$ são constantes positivas.
\begin{itemize}
\item[a)] Encontre a frequência e a amplitude de oscilações quando $w_0=1$, $\alpha=.1$ e $A=10$. (Teste diversas condições iniciais)
\item[b)] Estude a dependência da frequência e da amplitude com os parâmetros  $A$, $\alpha$ e $w_0$. (Teste diversas condições iniciais)
\item[c)] Que diferenças existem entre esse oscilador não linear e o oscilador linear?
\end{itemize}
\end{exer}

\begin{exer} Considere o seguinte modelo para um oscilador não linear:
\begin{eqnarray}
u''(t)-\alpha(A-z(t))u'(t)+w_0^2 u(t)&=&0\\
Cz'(t)+z(t)&=&u(t)^2
\end{eqnarray}
onde $A$, $\alpha$, $w_0$ e $C$ são constantes positivas.
\begin{itemize}
\item[a)] Encontre a frequência e a amplitude de oscilações quando $w_0=1$, $\alpha=.1$, $A=10$ e $C=10$. (Teste diversas condições iniciais)
\item[b)] Estude a dependência da frequência e da amplitude com os parâmetros  $A$, $\alpha$, $w_0$ e $C$. (Teste diversas condições iniciais)
\end{itemize}
\end{exer}

\begin{exer} Considere o seguinte modelo para o controle de temperatura em um processo químico:
\begin{eqnarray}
CT'(t)+T(t)&=&\kappa P(t)+T_{ext}\\
P'(t)&=&\alpha(T_{set}-T(t))
\end{eqnarray}
onde $C$, $\alpha$ e $\kappa$ são constantes positivas e $P(t)$ indica a potência do aquecedor. Sabendo que $T_{set}$ é a temperatura desejada, interprete o funcionamento desse sistema de controle e faça o que se pede:
\begin{itemize}
\item[a)] Calcule a solução quando a temperatura externa $T_{ext}=0$, $T_{set}=1000$, $C=10$, $\kappa=.1$ e $\alpha=.1$. Considere condições iniciais nulas.
\item[b)] Quanto tempo demora o sistema para atingir a temperatura 900K?
\item[c)] Refaça os dois primeiros itens com $\alpha=0,2$ e $\alpha=1$
\item[d)] Faça testes para verificar a influência de $T_{ext}$, $\alpha$ e $\kappa$ na temperatura final.
\end{itemize}
\end{exer}

\begin{exer} Considere a equação do pêndulo dada por:
\begin{equation} \frac{d^2\theta(t)}{dt^2}+\frac{g}{l}\sin(\theta(t))=0 \end{equation}
onde $g$ é o módulo da aceleração da gravidade e $l$ é o comprimento da haste.
\begin{itemize}
\item[a)] Mostre analiticamente que a energia total do sistema dada por
\begin{equation} \frac{1}{2}\left(\frac{d\theta(t)}{dt}\right)^2-\frac{g}{l}\cos(\theta(t)) \end{equation}
é mantida constante.
\item[b)] Resolva numericamente esta equação para $g=9,8m/s^2$ e $l=1m$ e as seguintes condições iniciais:
\subitem i. $\theta(0)=0,5$ e $\theta'(0)=0$.
\subitem ii. $\theta(0)=1,0$ e $\theta'(0)=0$.
\subitem iii. $\theta(0)=1,5$ e $\theta'(0)=0$.
\subitem iv. $\theta(0)=2,0$ e $\theta'(0)=0$.
\subitem v. $\theta(0)=2,5$ e $\theta'(0)=0$.
\subitem vi. $\theta(0)=3,0$ e $\theta'(0)=0$.
\end{itemize}
Em todos os casos, verifique se o método numérico reproduz a lei de conservação de energia e calcule período e amplitude.
\end{exer}

\begin{exer} Considere o modelo simplificado de FitzHugh-Nagumo para o potencial elétrico sobre a membrana de um neurônio:
\begin{eqnarray}
\frac{d V}{dt}& = &  V-V^3/3 - W +  I  \\
\frac{d W}{dt} & = & 0,08(V+0,7 - 0,8W)
\end{eqnarray}
onde $I$ é a corrente de excitação.
\begin{itemize}
\item Encontre o único estado estacionário $\left(V_0,W_0\right)$ com $I=0$.
\item Resolva numericamente o sistema com condições iniciais dadas por $\left(V_0,W_0\right)$ e
\subitem $I=0$
\subitem $I=0,2$
\subitem $I=0,4$
\subitem $I=0,8$
\subitem $I=e^{-t/200}$
\end{itemize}
\end{exer}



%\end{document}
